% concatenated from arxiv_version.tex
\documentclass[11pt]{article}

% Page layout (optional but recommended for arXiv)
\usepackage[margin=1in]{geometry}

% Recommended packages
\usepackage{microtype}
\usepackage{graphicx}
\usepackage{subcaption}
\usepackage{booktabs}
\usepackage{bm}
\usepackage{hyperref}
\usepackage[numbers]{natbib}
\usepackage{amsmath}
\usepackage{amssymb}
\usepackage{mathtools}
\usepackage{amsthm}
\usepackage{thmtools}
\usepackage[shortlabels]{enumitem}
\usepackage[capitalize,noabbrev]{cleveref}
\usepackage{algorithm}
\usepackage{algorithmic}
\usepackage{float}   % needed for [H]
%%%%%%%%%%%%%%%%%%%%%%%%%%%%%%%%
% THEOREMS
%%%%%%%%%%%%%%%%%%%%%%%%%%%%%%%%
\theoremstyle{plain}
\newtheorem{theorem}{Theorem}[section]
\newtheorem{proposition}[theorem]{Proposition}
\newtheorem{lemma}[theorem]{Lemma}
\newtheorem{fact}[theorem]{Fact}
\newtheorem{corollary}[theorem]{Corollary}
\theoremstyle{definition}
\newtheorem{definition}[theorem]{Definition}
\newtheorem{assumption}[theorem]{Assumption}
\theoremstyle{remark}
\newtheorem{remark}[theorem]{Remark}

% Todonotes (optional)
%\usepackage[disable,textsize=tiny]{todonotes}
\setlength{\marginparwidth}{2cm}
\usepackage[textsize=tiny]{todonotes}
\setlength{\marginparwidth}{4.5em}
\newcommand{\massi}[2][]{\todo[color=brown!20,#1]{{\bf MP:} #2}}
\newcommand{\matia}[2][]{\todo[color=red!20,#1]{{\bf MB:} #2}}
\newcommand{\saverio}[2][]{\todo[color=blue!20,#1]{{\bf SS:} #2}}

\newcommand{\MP}[1]{\textcolor{brown}{#1}}
\newcommand{\MB}[1]{\textcolor{red}{#1}}
\newcommand{\SaS}[1]{\textcolor{blue}{#1}}

%%%%%%%%%%%%%%%%%%%%%%%%%%%%%%%%
% MACROS
%%%%%%%%%%%%%%%%%%%%%%%%%%%%%%%%
\newcommand{\norm}[1]{\lVert{#1}\rVert}
\newcommand{\abs}[1]{\lvert{#1}\rvert}
\newcommand{\scalarp}[1]{\langle{#1}\rangle}
\newcommand{\pair}[2]{{\langle{{#1},{#2}}\rangle}}
\DeclareMathOperator*{\argmin}{argmin}
\DeclareMathOperator*{\argmax}{\mathrm{argmax}}
\DeclareMathOperator{\prox}{prox}
\newcommand{\HH}{X}
\newcommand{\bHH}{\bm{X}}
\newcommand{\R}{\mathbb{R}}
\newcommand{\Sb}{\mathbb{S}}
\newcommand{\x}{x}
\newcommand{\m}{d}
\newcommand{\bx}{\bm{x}}
\newcommand{\by}{\bm{y}}
\newcommand{\y}{y}
\newcommand{\buu}{\bm{g}}
\newcommand{\uu}{g}
\newcommand{\w}{w}
\newcommand{\iter}{n}
\newcommand{\ite}{k}
\newcommand{\N}{\mathbb{N}}
\newcommand{\RX}{\left]-\infty, +\infty\right]}
\newcommand{\reg}{\varphi}
\newcommand{\Id}{\mathrm{Id}}
\newcommand{\W}{W}
\newcommand{\balpha}{\boldsymbol{\alpha}}

%%%%%%%%%%%%%%%%%%%%%%%%%%%%%%%%
% TITLE
%%%%%%%%%%%%%%%%%%%%%%%%%%%%%%%%
\title{AdaGrad-Diff: A New Version of the Adaptive Gradient Algorithm}

\author{
Matia Bojovi\'c$^{1,2}$ \and
Saverio Salzo$^{1,3}$ \and 
Massimiliano Pontil$^{1,4}$
}


% \date{}

\date{} % remove date

\begin{document}

\maketitle

\footnotetext[1]{
Computational Statistics and Machine Learning, Istituto
Italiano di Tecnologia, Genoa, Italy}
\footnotetext[2]{
Department of Mathematics, University of Genoa, Genoa, Italy}
\footnotetext[3]{DIAG, Sapienza University of Rome, Rome, Italy}
\footnotetext[4]{Department of Computer Science, University College London, London, UK}

\begin{abstract}
%\massi{Ho messo una variazione del titolo. Un'altra possibilita' e' tenere solo "A New Version..." e dire AG-Diff nella'abstract}  
%\matia{Secondo me può essere  conveniente tenere AdaGrad nel titolo, dato che è una key word molto ricercata}
\noindent
Vanilla gradient methods are often highly sensitive to the choice of stepsize, which typically requires manual tuning. Adaptive methods alleviate this issue and have therefore become widely used. Among them, AdaGrad has been particularly influential. In this paper, we propose an AdaGrad-style adaptive method in which the adaptation is driven by the cumulative squared norms of successive gradient differences rather than gradient norms themselves. The key idea is that when gradients vary little across iterations, the stepsize is not unnecessarily reduced, while significant gradient fluctuations, reflecting curvature or instability, lead to automatic stepsize damping. Numerical experiments demonstrate that the proposed method is more robust than AdaGrad in several practically relevant settings.

% Gradient Descent and its stochastic variants remain among the most widely used first-order methods in machine learning. A central limitation of vanilla gradient methods, however, is that their update rules are agnostic to the geometry of the problem, leaving a highly sensitive stepsize to be manually tuned. Adaptive methods have become extremely popular precisely because they ease this burden. Among these, AdaGrad has had a profound influence, inspiring a large family of adaptive algorithms.
% We propose an AdaGrad-style adaptive variant, where the adaptation is driven not by the cumulative
% gradient squared norms, but by the cumulative squared norms of successive gradient differences. The key idea is that if gradients remain stable across iterations, the method does not shrink the stepsize unnecessarily. However, when gradients fluctuate significantly, indicating curvature, noise, or instability, the accumulated differences grow and the method automatically dampens the stepsize. In this way, the proposed method focuses on stabilizing learning in volatile directions while retaining efficiency in smooth, predictable ones. Preliminary numerical experiments indicate that the proposed variant is a viable alternative to AdaGrad in the scenarios we examined.
\end{abstract}

\section{Introduction} 

Gradient-based optimization methods are among the most widely used tools in machine learning due to their simplicity and broad applicability. In particular, Gradient Descent (GD) and its stochastic variants remain the canonical first-order methods for large-scale optimization problems. GD generates iterates according to
$\bx^{\iter+1} = \bx^\iter - \eta_\iter \buu^\iter$, starting from an arbitrary initial point $\bx^1$, where $\buu^\iter$ denotes a (sub)gradient of the objective function evaluated at $\bx^\iter$ and $\eta_\iter$ is the stepsize. The stepsize $\eta_\iter$ is the key parameter governing convergence speed and stability. Its choice is notoriously delicate: too small a value leads to slow progress, while too large a value may prevent convergence altogether. \\

\noindent
Classic convergence analysis of the GD algorithm for convex nonsmooth functions relies on conditions on 
the positive stepsizes $ \eta_\iter $ \cite{robbins&monro}. 
In particular, sufficient conditions are that $ (\eta_\iter)_{\iter\in \N}$ is a deterministic sequence 
of non-negative numbers that are square summable but whose sum diverges.  %satisfies
%\begin{equation*}   \sum_{\iter=1}^{\infty} \eta_\iter = \infty    \qquad\text{and}\qquad   \sum_{\iter=1}^{\infty} \eta_\iter^2 < \infty.
%\end{equation*}
In the finite-horizon setting, \citet{nesterov2018lectures} showed that the choice 
\begin{equation}\label{optimal_stepsize_Nesterov}
    (\eta_\ite)_{1\leq \ite\leq \iter} \equiv \frac{D}{G\sqrt{\iter}}
\end{equation}
is worst-case optimal in the complexity class, where $D := \lVert \bx^1 - \bx^\star \rVert$ being the starting distance from a minimizer $\bx^\star$, $G$ is the Lipschitz constant and $\iter$ is the number of total iterations. 
Similarly, in the convex $L$-Lipschitz smooth case, a sufficient condition for the convergence of GD is that the stepsizes satisfy
\begin{equation}\label{constraints_Lsmooth_analysis}
    \sup_{n\in\N}\eta_n < \frac{2}{L}.
\end{equation}
However, the quantities appearing in \eqref{optimal_stepsize_Nesterov} and \eqref{constraints_Lsmooth_analysis} are typically unknown in practice, making principled stepsize selection difficult and often requiring extensive hyperparameter tuning. \\

\noindent
To overcome these limitations, a broad body of work has focused on adaptive gradient methods, which automatically adjust stepsizes based on the observed optimization history. Among these methods, AdaGrad \cite{duchi2011adaptive} has played a particularly influential role, giving rise to a large family of adaptive algorithms. Its central idea is to replace a fixed or pre-scheduled stepsize by a sequence obtained by rescaling a base parameter $\eta>0$ using accumulated information from past gradients. One of the defining features of AdaGrad is that the stepsize adaptation is performed at the level of individual components. In its most general form, the iterates satisfy
\begin{equation}\label{adagrad_uncostrained_update}
x_i^{\iter+1} = x_i^\iter - \frac{\eta}{\w_i^\iter}\,\uu_i^\iter,
\end{equation}
where each $\w_i^\iter$ is a nondecreasing quantity that cumulates the squares of the $i$-th components of the gradients along the optimization trajectory. As $\w_i^\iter$ grows, the corresponding scaling factor $\eta/\w_i^\iter$ decreases, automatically moderating progress in directions where gradients accumulate rapidly. This component-wise adaptation enables AdaGrad to adjust to the geometry of the problem, alleviating the need to pre-specify quantities such as the Lipschitz constants $G$ and $L$ introduced above.\\

\noindent
Originally, convergence analysis was established in the $G$-Lipschitz continuous setting, and was later extended to the $L$-Lipschitz smooth case. Remarkably, in this extension, there are no restrictions on the choice of the stepsize parameter $\eta$ appearing in \eqref{adagrad_uncostrained_update}, in contrast to the constraints in \eqref{constraints_Lsmooth_analysis}. However, although convergence of the method is guaranteed for any $\eta > 0$, this does not imply that all choices yield comparable performance in practice. The scalar $\eta$ therefore remains a hyperparameter requiring tuning, and poorly chosen values can significantly degrade optimization performance.\\

\noindent
In this paper, we depart from the classical AdaGrad framework and introduce
a new adaptive policy, based on a difference-based mechanism, for a proximal gradient algorithm. We address 
%aimed at reducing sensitivity to the choice of the stepsize parameter $\eta$, while retaining the efficiency of the optimization. We address  
composite optimization problems of the form
\begin{equation}
    \min_{\bx\in \bHH}  F(\bx) := f(\bx) + \reg(\bx), 
\label{eq:main_problem}
\end{equation}
where $\bHH = \bigoplus_{i=1}^\m\HH_i$ is the direct sum of $\m\geq 1$ Hilbert spaces
$\HH_1, \dots, \HH_\m$, $f\colon \bHH\to \R$ is a convex function
and $\reg\colon \bHH\to \RX$ is an extended real-valued convex and lower semicontinuous function.
A special case of \eqref{eq:main_problem} is that of $\bHH= \R^\m = \R\oplus \cdots \oplus \R$ ($\m$ times), endowed with the standard Euclidean scalar product, $f$ representing a loss function and $\reg$ an appropriate regularizer.
%In particular \matia{Da qua in poi lo metterei nella sezione esperimenti} the loss function $f$ may have the form $f(\x) = \frac{1}{N}\sum_{j=1}^N \ell(\x, z_i)$, where $z_i = (a_i, b_i)$ represents an input-output pair drawn from an (unknown) underlying distribution $\mathcal{D}$. Some possibilities for $\ell$ include
%\begin{itemize}   \item \textbf{Least-squares:} $\ell(\x, (a, b)) = (b - \scalarp{\x, a})^2$, with $a \in \R^n$, and $b \in \R$.     \item \textbf{Hinge loss:} $\ell(\x, (a, b)) = \max(0, 1 - b(\scalarp{\x, b}))$, with $a \in \R^n$ and $b \in \{+1, -1\}$. \item \textbf{Logistic regression:} $\ell(\x, (a, b)) = \log(1 + \exp(-b(\scalarp{a,x})))$, with $x \in \R^n$, $y \in \{+1, -1\}$. \end{itemize}
%Examples of the regularization term $\reg(\x)$ include:
%\begin{itemize}   \item \textbf{$\ell_1$-regularization:} $\reg(\x) = \lambda \|\x\|_1$ with $\lambda > 0$. \item \textbf{$\ell_2$-regularization:} $\reg(\x) = (\sigma/2) \|\x\|_2^2$ with $\sigma > 0$. \item \textbf{Convex constraints:} $\reg(\x)=\iota_C(\x)$ is the \textit{indicator function} of a closed convex set $C$, i.e., $\reg(\x) = 0$ if $\x \in C$ and $+\infty$ otherwise. \end{itemize}

\subsection{Related works}
Gradient-based optimization methods form the foundation of machine learning and signal processing. Among these approaches, adaptive methods have proven especially effective in practice, with AdaGrad serving as a foundational algorithm that has inspired several extensions. The original AdaGrad algorithm, designed for stochastic online optimization, was proposed simultaneously by \cite{duchi2011adaptive} and \cite{streeter2010less}.\\

\noindent
Subsequently, a number of modifications of AdaGrad have been proposed, particularly with the aim of avoiding the continual decay of the stepsizes and the need of tuning the parameter $\eta$.
In this line of works, RMSProp \cite{hinton} replace AdaGrad’s cumulative accumulation of squared gradients with exponential moving averages to prevent excessively rapid stepsize decay.
AdaDelta \cite{zeiler2012adadelta} improves RMSProp by removing the parameter $\eta$, by leveraging approximate Hessian information.  Next, Adam \cite{kingma2014adam} further combines RMSProp-like aggregation of the gradients with momentum, yielding an optimizer that has become widely used and effective in practice. 
In the spirit of AdaDelta, which seeks to remove the need for a manually chosen stepsize $\eta$, several works have developed parameter-free algorithms based on the AdaGrad framework. Building on the theoretically optimal stepsize in \eqref{optimal_stepsize_Nesterov}, \citet{ivgi2023dog}, and \citet{defazio2023learning} introduce additional adaptivity in the numerator of the stepsize, namely setting $\eta_\iter = \gamma_\iter/\w^\iter$, where the sequence $(\gamma_\iter)_{\iter\in\N}$ serves as an estimate of the constant $D$. 
Other methods achieve last-iterate convergence rates in the $L$-Lipschitz smooth convex setting by appropriately reweighting the cumulative sum of the squared gradients \cite{liu2022convergence}. 
Finally, certain approaches tackle undesirable scale-dependence by removing the square-root in the AdaGrad update \cite{choudhury2024remove}, extending applicability to $L$-Lipschitz smooth, potentially nonconvex objectives. These contributions highlight that, 
carefully designed variants of the stepsize updating rule in AdaGrad may substantially improve performances in real applications.\\

\noindent
Our contribution belongs to this stream of works. Indeed, we propose a new modification of AdaGrad’s adaptivity rule that is motivated by stability considerations arising in practice, while simultaneously providing theoretical guarantees. More specifically, instead of accumulating sums of squared gradients, we compute cumulative sums of squared gradient differences. This proposal does not alter the underlying update structure and retains AdaGrad’s practical appeal, while addressing limitations identified in both empirical and theoretical settings.\\

\noindent
Of particular relevance to the present work are two recent results on vanilla AdaGrad for $L$-Lipschitz smooth convex objectives in a finite-dimensional setting. \citet{traore2021sequential} establish, for the first time, the convergence of the iterate. On the other hand, \citet{liu2022convergence} derive explicit convergence rates in the objective gap, although with an exponential dependence on the dimension of the space.\\

\noindent
We conclude this section by briefly discussing the stochastic setting.
One main difficulty in handling  stochastic gradients in AdaGrad-like schemes 
is that the stepsizes become stochastic random variables which appear multiplicatively in front of the gradients. To overcome this issue appropriate modifications of the stepsize updating rule have been proposed. \citet{li2019convergence} remove the most recent gradient from the cumulative sum of squared gradients, ensuring that $\eta_{\iter}$ is conditionally independent on the current gradient $\buu^\iter$.
On the other hand, \citet{ward2020adagrad} analyze the vanilla AdaGrad scheme and introduce, in their analysis, a proxy stepsize $\tilde{\eta}_\iter$ that is conditionally independent on $\buu^\iter$. This construction allows them to decouple the gradient and the stepsize. 
Building on the same technique, \citet{defossez2020simple} employ a similar argument to establish theoretical guarantees for Adam.
Finally, \citet{faw2022power} further generalize the results of \citet{ward2020adagrad} by relaxing the assumptions on the gradient noise.

\subsection{Contributions}
% We propose a new algorithm based on the AdaGrad step-size rule with a single key innovation: we accumulate the norms of successive gradient differences instead of gradient norms. The intuition is straightforward: when the function’s landscape remains relatively stable along the optimization trajectory, that is, when the gradient does not fluctuate significantly, the method maintains larger stepsizes without unnecessary penalty. Conversely, when gradients change significantly between iterations, the accumulated differences trigger stepsize reduction to stabilize the optimization dynamics. 
In this section, we collect the main contributions. 
\begin{itemize}
\item We propose a new algorithm based on a modification of the AdaGrad stepsize rule with a single key innovation: we accumulate the norms of successive gradient differences instead of gradient norms. 
\item We provide convergence rates in the objective gap for problem \eqref{eq:main_problem} when $f$ is either $G$-Lipschitz continuous or $L$-Lipschitz smooth, retaining $\mathcal{O}\left(1/\sqrt{n}\right)$ or $\mathcal{O}\left(1/n\right)$ decay, respectively.
\item We provide (weak) convergence of the iterates to a minimizer for problem \eqref{eq:main_problem} when $f$ is $L$-Lipschitz smooth. To the best of our knowledge, in the composite setting, such guarantee has not been established for AdaGrad, even for finite dimension.
\item Finally, we evaluate our algorithm on a range of problems. 
Experiments show that AdaGrad-Diff consistently performs well over a broader range of values of the parameter $\eta$ compared to AdaGrad. Moreover, the algorithm naturally mitigates suboptimal choices of $\eta$: when $\eta$ is too small, the effective stepsize adapts to allow more aggressive progress, while overly large values of $\eta$ are moderated through the same mechanism. Overall, this adaptive behavior enhances stability without sacrificing efficiency and substantially reduces the need for extensive hyperparameter tuning.
% In the cases considered, the method automatically reduces the stepsize in directions where gradients change rapidly between iterations, while maintaining larger stepsizes when the gradient remains stable. The experiments show that this behavior improves stability without sacrificing efficiency and makes the algorithm less sensitive to the choice of $\eta$, reducing the need for extensive hyperparameter tuning.
\end{itemize}

\subsection{Organization of the paper}

The remainder of this section introduces the notation used throughout the paper. In Section~\ref{2}, we present the algorithmic framework and state the main theoretical results. Section~\ref{3} develops the key preparatory lemmas and technical tools required for the analysis. Section~\ref{4} reports numerical experiments illustrating the performance of the proposed method. Some proofs and additional experimental results are deferred to the Appendix.
\subsection{Notation}

We denote by $\N$ the set of natural numbers excluding zero and set
$\R_+ = \left[0,+\infty\right[$ and $\R_{++} = \left]0,+\infty\right[$.
For every $\m\in \N$, we also set $[\m]= \{1,\dots, \m\}$. If $\HH$
is a Hilbert space we denote by $\scalarp{\cdot, \cdot}$ and $\norm{\cdot}$
the scalar product and norm of $\HH$ respectively. 
%We denote by $\Sb(\HH)$ the space of self-adjoint and bounded linear operators on $\HH$.
The Loewner ordering on linear bounded self-adjoint operators on $\HH$
 is defined as $\W_1 \preceq \W_2 \Leftrightarrow\ (\forall\,\x\in \HH)\ \scalarp{\W_1\x,\x} \leq \scalarp{\W_2\x,\x}$. When $\W\succeq 0$ we say that $\W$ is positive.
We denote by $\Gamma_0(\HH)$ the set of extended real-valued functions
$f\colon \HH\to \RX$ which are proper convex and lower semicontinuous.
If $f\in \Gamma_0(\HH)$ its subdifferential is the set-valued mapping
$\partial f\colon \HH\to 2^\HH$ such that, for every $\x\in \HH$, $\partial f(\x)=\{\uu\in \HH\,\vert\, \forall\,\y\in \HH\colon f(\y)\geq f(\x) + \scalarp{\y-\x, \uu}\}$.
If $\HH_1, \dots, \HH_{\m}$
are Hilbert spaces their direct sum is the vector space $\bHH= \prod_{i=1}^\m \HH_i$  
endowed with the scalar product $\scalarp{\bx, \by} = \sum_{i=1}^\m \scalarp{\x_i, \y_i}$, where
$\bx=(\x_i)_{1\leq i\leq \m}$ and $\by=(\y_i)_{1\leq i\leq \m}$. This Hilbert space is denoted by $\bigoplus_{i=1}^\m \HH_i$. If $\w_1, \cdots, \w_m\in \R_{++}$ and $\Id_i$ denote the identity operator of $\HH_i$ we denote by $\W = \bigoplus_{i=1}^\m \w_i \Id_i\colon \bHH\to \bHH$ the block-diagonal operator $\bx \mapsto (\w_i \x_i)_{1\leq i\leq \m}$, which is positive self-adjoint bounded operator with $\alpha \Id \preceq \W_\iter\preceq \beta\Id$, where $0<\alpha:= \min_{1\leq i\leq \m}\w_i \leq \max_{1\leq i\leq \m}\w_i=\norm{\W_n}=:\beta$. Thus, we define the scalar product on $\bHH= \bigoplus_{i=1}^\m \HH_i$, as
$\scalarp{\bx,\by}_{\W}= \scalarp{\W \bx, \by}=\sum_{i=1}^\m \w_i \scalarp{x_i, y_i}$. 
%The corresponding norm is $\norm{\bx}_{\W}= (\scalarp{\W \bx, \bx})^{1/2} = (\sum_{i=1}^\m w_i \norm{x_i}^2)^{1/2}$.


\section{Algorithm and Main Results}\label{2}

In this section we make precise the setting, we present the algorithm, and finally the main results of the paper.

\begin{assumption}\label{ass:0}
$\bHH$ is a Hilbert space which is the direct sum of $\m\geq 1$ Hilbert spaces $\HH_1, \dots, \HH_\m$.
The function $f\colon \bHH\to \R$ is convex and the function $g\in\Gamma_0(\bHH)$. Moreover, $F := f+g$ and $\argmin F\neq \varnothing$. We denote the minimal value of $F$ by $F_\star$.
\end{assumption}
\noindent
We consider two additional alternative assumptions for $f$.\\
\begin{assumption}\label{ass:1}
$f$ is $G$-Lipschitz continuous with $G>0$, i.e.,
\begin{equation*}
    \forall\,\bx,\by\in \bHH\colon \abs{f(\bx)-f(\by)}\leq G\norm{\bx-\by}.
\end{equation*}
\end{assumption}

\begin{assumption}\label{ass:2}
$f$ is $L$-Lipschitz smooth with $L>0$, i.e., it is differentiable and
\begin{equation*}
    \forall\,\bx,\by\in \bHH\colon \norm{\nabla f(\bx)-\nabla f(\by)}\leq L\norm{\bx-\by}.
\end{equation*}
\end{assumption}
\noindent
We now introduce the algorithm.
We consider a variable metric setting defined by a sequence $(\W_\iter)_{\iter\in \N}$ of positive definite block-diagonal bounded linear operators
\begin{equation*}
\W_\iter\colon \bHH\to \bHH\colon \bx\mapsto 
%(\w^{(\iter)}_i x_i)_{1\leq i\leq \m},
\begin{bmatrix}
\w_1^n \x^1\\
\vdots\\
\w_\m^n \x^\m\\
\end{bmatrix}
\end{equation*}
where $\w_i^\iter>0$ are properly set.
For every $\iter\in \N$, $\W_\iter$ defines the norm
\begin{equation}
    \|\bx\|_{\W_\iter} := \sqrt{\scalarp{\W_\iter \bx, \bx}} = \sqrt{\sum_{i=1}^\m \w_i^{\iter} \norm{x_i}^2}.
\end{equation}
\noindent
Then, for the current iteration $\bx^\iter$, given a parameter $\eta_\iter > 0$ and 
 $\buu^\iter\in \partial f(\bx^\iter)$, the \emph{proximal (sub)gradient} update in the geometry induced by $\W_\iter$ is defined as
 \begin{align}
    \label{eq:prox-update-general}
    \nonumber\bx^{\iter+1} \!&= 
    \argmin_{\bx}\left\{\scalarp{\buu^\iter, \bx-\bx^\iter} + \reg(\bx) + \frac{\norm{\bx-\bx^\iter}_{\W_\iter}^2}{2\eta_\iter} \right\}\\
    &=\prox^{\W_n}_{\eta_\iter \reg}(\bx^\iter - \eta_\iter \W_\iter^{-1}\buu^\iter).
\end{align}
A classical choice for the diagonal geometry is the \emph{AdaGrad metric}
\begin{equation*}
%\label{AdaGrad_stepsizes}
(\forall\,i\in [\m])\  \   \w_{i}^\iter
    := \varepsilon+\sqrt{\sum_{\ite=1}^\iter \norm{\uu^\ite_{i}}^2},\ \ \eta_\iter \equiv \eta,\ \varepsilon > 0,
\end{equation*}
where, for every $\ite\in \N$, $\buu^\ite = (\uu_i^\ite)_{1\leq i\leq \m}$.
When $\reg \equiv 0$, the proximal gradient step above reduces to the standard (unconstrained) AdaGrad update, namely
\begin{equation}
(\forall\,i\in [\m])\ \    \x_i^{\iter+1} = \x_{i}^\iter - \frac{\eta}{\varepsilon +   \sqrt{\sum_{\ite=1}^\iter \norm{\uu_{i}^\ite}^2}} \uu_{i}^\iter.
\end{equation}
For general $\reg$, the update \eqref{eq:prox-update-general} combines adaptive geometry with regularization.\\

\noindent
Our method follows the composite and geometry-aware structure of \eqref{eq:prox-update-general}. The key departure from AdaGrad is in the construction of the metric $\W_\iter$: rather than accumulating squared subgradient norms, we accumulate the squared norms of successive subgradient differences. 
Specifically, the positive definite block-diagonal operator $\W_\iter$ is defined by the following parameters
\begin{equation}\label{AdaDiff}
(\forall\,i\in [\m])\ \ \w_{i}^{\iter} := \varepsilon+\sqrt{\sum_{\ite=1}^\iter \norm{\uu_{i}^\ite - \uu_{i}^{\ite-1}}^2}\;, 
\end{equation}
\noindent
where we use the convention that $\buu^0 = 0$. This construction directly encodes our key insight: when subgradients remain stable across iterations, the accumulated differences grow slowly and the stepsize remains large; on the contrary when subgradients fluctuate significantly, the accumulated differences grow rapidly and the stepsize reduces to stabilize the optimization. 
The resulting algorithm maintains the composite structure and diagonal geometry of AdaGrad while providing refined adaptation that distinguishes between stable and volatile directions during optimization. 

\begin{algorithm}[H]
\caption{AdaGrad-Diff}\label{adadiff}
\begin{algorithmic}[1]
\STATE \textbf{Input:} $\bx^1 \in \bHH$, $\eta > 0$ and $\varepsilon > 0$.
\STATE Set $\buu^0 = 0\in \HH$. 
\FOR{$\iter = 1,2, \dots,$}
    \STATE Let $\buu^\iter \in \partial f(\bx^\iter)$
    \STATE $(\forall\, i\in [\m])\ \w^\iter_i = \varepsilon+\left(\sum_{\ite=1}^\iter \norm{\uu_i^\ite - \uu_i^{\ite-1}}^2\right)^{1/2}$  
    \STATE $\W_\iter = \bigoplus_{i=1}^\m \w_i^\iter \Id_i\phantom{\Big(}$
    \STATE $\bx^{\iter+1} = \prox^{\W_\iter}_{\eta\reg} (\bx^\iter - \eta \W_\iter^{-1} \buu^\iter )$
\ENDFOR
\end{algorithmic}
\end{algorithm}
\noindent
We now present the main results
\begin{theorem}\label{thm_nonsmooth}
Under Assumptions~\ref{ass:0} and \ref{ass:1}, let $(\bx^\iter)_{n\in\N}$ be the sequence generated by Algorithm~\ref{adadiff} and suppose that $(\bx^\iter)_{n\in\N}$ is bounded.\footnote{This can be ensured if $\mathrm{dom}\,\reg$ is bounded.} Define, for every $n\in \N$, $\bar{\bx}^\iter = (\iter-1)^{-1}\sum_{\ite=2}^\iter \bx^\ite$. Then, the following rate of convergence hold
    \begin{equation}
        F(\bar{\bx}^\iter) - F_\star = \mathcal{O}\left(\frac{1}{\sqrt{\iter}}\right).
    \end{equation}
\end{theorem}
\begin{theorem}\label{thm_smooth}
Under Assumptions~\ref{ass:0} and \ref{ass:2}, let $(\bx^\iter)_{n\in\N}$ be the sequence generated by Algorithm~\ref{adadiff} and define, for every $n\in \N$, $\bar{\bx}^\iter = (\iter-1)^{-1}\sum_{\ite=2}^\iter \bx^\ite$. Then, the following rate of convergence hold
    \begin{equation}
    \label{eq:rate2}
        F(\bar{\bx}^\iter) - F_\star = \mathcal{O}\left(\frac{1}{\iter}\right).
    \end{equation}
Moreover, $\bx^n \rightharpoonup \bx^*$ for some $\bx^*\in \argmin F$.
\end{theorem}






\section{Convergence Analysis}\label{3}
\label{sec:analysis}
In this section we provide the main arguments used for proving the two theorems showed in the previous section. All the proofs can be found in Appendix \ref{A}. 
In what follows we assume that the standing Assumption~\ref{ass:0} holds and that the sequences
 $(\bx^\iter)_{n\in\N}$ and $(\buu^\iter)_{n\in \N}$ are generated according to the update in \eqref{eq:prox-update-general} with a general sequence of stepsizes $(\eta_\iter)_{\iter\in\N}$.\\

\noindent
We start with a fundamental result concerning two basic inequalities related to the proximal gradient step.


\begin{lemma}[{\bf The basic inequalities}]\label{lemma1}
%Under Assumption~\ref{ass:0}, 
%let $(\x^\iter)_{\iter\in \N}$ be the iterates generated Algorithm~\ref{adadiff}, 
Let $\bx\in \bHH$ and $\iter \in \N$. Then 
\begin{equation}\label{lemma1_bound}
    2\eta_\iter(F(\bx^{\iter+1}) - F(\bx))  
    \le \norm{\bx^\iter - \bx}_{\W_\iter}^2- \norm{\bx^{\iter+1} - \bx}_{\W_\iter}^2 + \eta_\iter^2 \norm{\buu^{\iter+1} - \buu^{\iter}}_{\W_\iter^{-1}}^2.
    \end{equation}
Moreover, under Assumption~\ref{ass:2}, it holds
\begin{align}\label{lemma1_bound2}
    \nonumber 2\eta_\iter(F(\bx^{\iter+1}) - F(\bx))  
    &\le \norm{\bx^\iter - \bx}_{\W_\iter}^2
    - \norm{\bx^{\iter+1} - \bx}_{\W_\iter}^2 \\
    &\qquad + \eta_\iter^2 
    \norm{\buu^{\iter+1} - \buu^{\iter}}_{\W_\iter^{-1}}^2 
    - \frac{\eta_\iter}{L}
    \norm{\buu^{\iter+1} - \buu^{\iter}}^2 .
\end{align}
\end{lemma}
\noindent
A comment is in order. There is a key departure from the classical bound provided for AdaGrad. Indeed in \cite{duchi2011adaptive} the following inequality is proved
\begin{align}\label{AdaGradBound}
    \nonumber 2\eta_\iter(F(\bx^{\iter}) - F(\bx))  & \le \lVert \bx^\iter - \bx \rVert_{\W_\iter}^2- \lVert \bx^{\iter+1} - \bx\rVert_{\W_\iter}^2  \\ & \qquad + \eta_\iter^2 \lVert \buu^{\iter} \rVert_{\W_\iter^{-1}}^2 + 2\eta_\iter(\reg(\bx^{\iter}) - \reg(\bx^{\iter+1})).
\end{align}
Our bound, on the other hand, involves the norm of successive gradient differences. \\
As a consequence of Lemma~\ref{lemma1}, we get the following Corollary.  

\begin{corollary}
\label{Proposition1}
Let $\bx\in \bHH$ and $\iter \in \N$ and set
%Under the same assumptions of Lemma~\ref{lemma1}, set 
$\Delta^{\iter-1}_{\max} = \max_{1\leq \ite\leq \iter-1} \max_{1\leq i\leq \m} \norm{\x^\ite_i - x_i}^2$. Then
\begin{align}\label{Proposition1_bound}
    \nonumber\sum_{\ite=2}^\iter\eta_{\ite-1}(F(\bx^\ite) - F(\bx)) &\le \frac{\varepsilon}{2}\lVert \bx^1 - \bx\rVert^2 + \frac{1}{2}\Delta_{\max}^{\iter-1} \sum_{i=1}^\m \w_i^{\iter-1} \\
     &\qquad  + \frac{1}{2}\sum_{\ite=2}^\iter\eta_{\ite-1}^2 \lVert \buu^{\ite} - \buu^{\ite-1} \rVert_{\W_{\ite-1}^{-1}}^2. 
    \end{align}
\end{corollary}
\noindent
From now on, we restrict our attention to Algorithm \ref{adadiff}.
\begin{proposition}\label{main_theorem}
Let $(\bx^\iter)_{\iter\in \N}$ be generated by Algorithm~\ref{adadiff}.
%Under Assumption~\ref{ass:0}, 
%let $(\x^\iter)_{\iter\in \N}$ be the iterates generated Algorithm~\ref{adadiff}, 
%Let $\x\in \HH$. Then for every $\iter \in \N$
Let $\bx\in \bHH$, $\iter \in \N$ and set
$\Delta^{\iter-1}_{\max} = \max_{1\leq \ite\leq \iter-1} \max_{1\leq i\leq \m} \norm{\x^\ite_i - x_i}^2$. 
Then
\begin{align}
\nonumber F(\bar{\bx}^{\iter}) - F(\bx)  &\le \frac{\varepsilon}{2\eta(\iter-1)}\lVert \bx^1 - \bx\rVert^2 +\frac{\sqrt{\m}}{\iter-1}\left[\frac{1}{2\eta}\Delta^{\iter-1}_{\max} + \eta \right]
\sqrt{\sum_{\ite=1}^\iter \norm{\buu^{\ite-1}- \buu^\ite}^2}\\
& \qquad + \frac{\eta}{2(\iter-1)} \sum_{i=1}^\m\sum_{\ite=2}^{\iter} \left(\frac{1}{\w_i^{\ite-1}} - \frac{1}{\w_i^\ite}\right)\norm{\uu_i^\ite - \uu_i^{\ite-1}}^2. \label{extra}
\end{align}
\end{proposition}


\subsection{The $G$-Lipschitz continuous case}
\label{sec:analysis_a}
In this section, beyond Assumptions~\ref{ass:0}, we make also Assumptions~\ref{ass:1}.
% Moreover, we take  the sequence of stepsizes
% constant, say  $\eta_\iter\equiv\eta>0$.
\begin{proof}[Proof of Theorem~\ref{thm_nonsmooth}]
Considering Proposition~\ref{main_theorem}, we need just to control the two sums in bound \eqref{extra}. 
This is possible since the assumptions ensure that the subgradients of $f$ are bounded in norm by $G$.
More specifically, as for the first sum we have
\begin{align*}
%\label{eq:20260127a}
\sqrt{ \sum_{\ite=1}^\iter \norm{\buu^\ite - \buu^{\ite-1}}^2} \leq 2 G \sqrt{n}.
\end{align*}
On the other hand for the second sum in \eqref{extra} we have
\begin{align*}
\sum_{i=1}^\m\sum_{\ite=2}^{\iter} \left(\frac{1}{\w_i^{\ite-1}} - \frac{1}{\w_i^\ite}\right)\norm{\uu_i^\ite - \uu_i^{\ite-1}}^2& \leq 4 G^2 \sum_{i=1}^\m
\sum_{\ite=2}^{\iter} \left(\frac{1}{\w_i^{\ite-1}} - \frac{1}{\w_i^\ite}\right) \leq 4 G^2 \sum_{i=1}^\m \frac{1}{\w_i^1}.
\end{align*}
In the end, we have
\begin{align*}
F(\bar{\bx}^{\iter}) - F(\bx) & \leq \frac{\varepsilon}{2\eta(\iter-1)}\lVert \bx^1 - \bx\rVert^2 +  \frac{2 G\sqrt{n d}}{\iter-1} \left[\frac{1}{2\eta}\Delta^{\iter-1}_{\max} + \eta \right] +  \frac{\eta 4 G^2}{2(\iter-1)}  \sum_{i=1}^\m \frac{1}{\w_i^1}
\end{align*}
and the statement follows by choosing $\bx\in\argmin F$, considering that $(\Delta^{\iter-1}_{\max})_{\iter\in \N}$ is also bounded by assumption.
\end{proof}


\subsection{The $L$-Lipschitz smooth case}
\label{sec:analysis_b}
Here we make Assumptions~\ref{ass:0} and Assumptions~\ref{ass:2}. \\
The bound in Proposition~\ref{main_theorem} shows that in order to get a rate of convergence 
for $F(\bar{\bx}^\iter)-F_*$, it is necessary to bound the sums $\sum_{\ite=1}^{\iter}\norm{\buu^\ite-\buu^{\ite-1}}^2$. On the other hand,
the convergence of the iterates for proximal gradient methods classically follows from the so-called Fej\'er monotonicity property.
See \cite{combettes2013variable} and the Appendix \ref{A} for details. In this respect,
inequality \eqref{lemma1_bound} in Lemma~\ref{lemma1} with $\bx\in \argmin F$ yields 
\begin{equation*}
    \norm{\bx^{\iter+1} - \bx}_{\W_\iter}^2\leq 
     \norm{\bx^\iter - \bx}_{\W_\iter}^2 + \eta^2 \norm{\buu^{\iter+1} - \buu^{\iter}}_{\W_\iter^{-1}}^2. 
\end{equation*}
This already resembles the Fej\'er monotonicity property in a variable metric setting,
except that on the left-hand side, it should appear the norm with respect to $\W_{\iter+1}$
and that again the sequence $(\norm{\buu^{\iter+1} - \buu^{\iter}}_{\W_\iter^{-1}}^2)_{\iter\in \N}$ should be summable. The following results show that indeed it is possible to address these issues.


\begin{proposition}\label{proposition5}
We have that $\sum_{\iter=1}^\infty \lVert \buu^{\iter+1} - \buu^\iter \rVert^2<+\infty$.
\end{proposition}
\noindent
The next result ensures that the sequence $(\bx^\iter)_{\iter\in \N}$ is 
quasi-Fej\'er  with respect to $\argmin F$ relative to $(\W_\iter)_{\iter\in \N}$.

\begin{proposition}
\label{prop:Fejer}
Under Assumptions~\ref{ass:0} and \ref{ass:2},
the sequence $(\bx^\iter)_{\iter\in \N}$ generated by Algorithm~\ref{adadiff} satisfies the following inequality  for every $\iter\in \N$ and $\bx\in \argmin F$
\begin{equation*}
\norm{\bx^{\iter+1}-\bx}^2_{\W_{\iter+1}} \leq (1+\chi_\iter) \norm{\bx^{\iter}-\bx}^2_{\W_{\iter}} + \alpha_\iter
\end{equation*}
where $(\chi_\iter)_{\iter\in \N}$ and $(\alpha_\iter)_{\iter\in \N}$ are summable, defined below
\begin{equation*}
\chi_\iter :=\max_{1\leq i\leq \m}\frac{\varepsilon + v_i^{\iter}}{\varepsilon + v_i^{\iter-1}}- 1, \quad
\alpha_\iter := \eta^2 \norm{\buu^{\iter+1} - \buu^{\iter}}_{\W_\iter^{-1}}^2.
\end{equation*}
\end{proposition}
\noindent
Once we have the Fej\'er monotonicity property, the convergence of the iterates is ensured by the following result. 
\begin{proposition}
\label{prop:clusterpoints}
Every weak sequential cluster point of $(\bx^{\iter})_{\iter\in \N}$ belongs to $\argmin F$.
\end{proposition}


\begin{figure*}[ht]
    \centering
    %\vskip 0.2in
    % Row 1
    \begin{subfigure}{0.374\textwidth}
        \includegraphics[width=\linewidth]{./hinge_synthetic.png}
    \end{subfigure}%
    \hspace{.5em}
    \begin{subfigure}{0.374\textwidth}
        \includegraphics[width=\linewidth]{./LAD_synthetic.png}
    \end{subfigure}
    \vspace{0.3em}  % less vertical space
    % Row 2
    \begin{subfigure}{0.374\textwidth}
        \includegraphics[width=\linewidth]{./tuned_losses_hinge.png}
        \caption{Hinge Loss on synthetic.}
    \end{subfigure}%
    \hspace{.5em}
    \begin{subfigure}{0.374\textwidth}
        \includegraphics[width=\linewidth]{./tuned_losses_LAD.png}
        \caption{LAD Regression on synthetic.}
    \end{subfigure}
    \caption{On the first row, final optimality gaps after $\iter$ iterations across different choices for the parameter $\eta$, illustrating the sensitivity of the two methods in the setting with nonsmooth losses. On the second row, the minimization performance of AdaGrad and AdaGrad-Diff with optimally tuned stepsizes. The plots show the average and standard deviation over 10 initialization of the algorithms.}
    \label{fig:robustness_nonsmooth}
\end{figure*}
\begin{proof}[Proof of Theorem~\ref{thm_smooth}]
We first address the rate of convergence \eqref{eq:rate2}.
Similarly to the proof of Theorem~\ref{thm_nonsmooth} we need to control the second and third terms in the bound \eqref{extra}. To that purpose it is critical to use the results in Proposition~\ref{proposition5}
and Proposition~\ref{prop:Fejer} which ensure that $\sum_{\iter=1}^\infty \lVert \buu^{\iter+1} - \buu^\iter \rVert^2<+\infty$ and that the sequence $(\bx^\iter)_{\iter\in \N}$, being quasi-Fej\'er, is bounded, so that
$(\Delta^{\iter-1}_{\max})_{\iter\in \N}$ is bounded too. Note that as for last term in \eqref{extra}, 
%considering that $1/\w_i^{\ite-1}\leq 1/\varepsilon$, 
we have
\begin{align*}
\sum_{i=1}^\m\sum_{\ite=2}^{\iter} &\left(\frac{1}{\w_i^{\ite-1}} - \frac{1}{\w_i^\ite}\right)\norm{\uu_i^\ite - \uu_i^{\ite-1}}^2 \leq \max_{1\leq \ite\leq \iter} \norm{\buu^{\ite}-\buu^{\ite-1}}^2\sum_{i=1}^\m \frac{1}{\w^1_i}.
%\leq \frac 1 \varepsilon \sum_{\ite=2}^{\iter}\norm{\buu^\ite - \buu^{\ite-1}}^2.
\end{align*}
Since the sequence $(\norm{\buu^{\iter}-\buu^{\iter-1}}^2)_{\iter \in \N}$ is summable, it is bounded and hence the right hand-side of the above inequality can be bounded by a constant.
Concerning the weak convergence of the iterates, in view of Fact~\ref{fact:fejer}, it is a direct consequence of Proposition~\ref{prop:Fejer} and Proposition~\ref{prop:clusterpoints}.
\end{proof}





%If you are using \LaTeX, please use the ``algorithm'' and ``algorithmic'' environments to format pseudocode. These require the corresponding stylefiles, algorithm.sty and algorithmic.sty, which are supplied with this package. \cref{alg:example} shows an example.

\begin{figure*}[ht]
    %\vskip 0.2in
    \centering
    
    % Row 1
    \begin{subfigure}{0.33\textwidth}
        \includegraphics[width=\linewidth]{./logistic_news20.png}
    \end{subfigure}%
    \begin{subfigure}{0.33\textwidth}
        \includegraphics[width=\linewidth]{./logistic_splice.png}
    \end{subfigure}%
    \begin{subfigure}{0.33\textwidth}
        \includegraphics[width=\linewidth]{./kernel_2moons.png}
    \end{subfigure}

    \vspace{0.3em}  % less vertical gap

    % Row 2
    \begin{subfigure}[t]{0.33\textwidth}
        \includegraphics[width=\linewidth]{./tuned_losses_news20.png}
        \caption{Log.~Regression on \texttt{news20}.}
    \end{subfigure}%
    \begin{subfigure}[t]{0.33\textwidth}
        \includegraphics[width=\linewidth]{./tuned_losses_splice.png}
        \caption{Log.~Regression on \texttt{splice.t}.}
    \end{subfigure}%
    \begin{subfigure}[t]{0.33\textwidth}
        \includegraphics[width=\linewidth]{./tuned_losses_2moons.png}
        \caption{SVM Classification on \texttt{2moons}.}
    \end{subfigure}

    \caption{
    On the first row, final optimality gaps after $\iter$ iterations across different choices for the parameter $\eta$, illustrating the sensitivity of the two methods in the setting with smooth losses. On the second row, the minimization performance of AdaGrad and AdaGrad-Diff with optimally tuned stepsizes. The plots show the average and standard deviation over 10 initializations of the algorithms.
    }
    \label{fig:robustness_smooth}
    \vspace{-.2truecm}
\end{figure*}

\begin{figure*}[ht]
    %\vskip 0.2in
    \centering
    \begin{subfigure}[t]{\textwidth}
        \centering
        \includegraphics[width=\linewidth]{./losses_splice.png}
        \caption{Optimality gap.}
        \label{fig:losses-splice}
    \end{subfigure}
    \hfill
    \begin{subfigure}[ht]{\textwidth}
        \centering
        \includegraphics[width=\linewidth]{./stepsizes_splice.png}
        \caption{Stepsizes evolution.}
        \label{fig:stepsizes-splice}
    \end{subfigure}
    \caption{Comparison of AdaGrad and AdaGrad-Diff for Logistic Regression on the \texttt{splice.t} dataset.
    (a) Optimality gaps  across three stepsizes parameter settings
    ($\eta = 0.0238,\, 0.238,\, 2.38$).
    (b) Stepsize evolution across different choices for the parameter $\eta$. The plots show the average and standard deviation over 10 initializations of the algorithms.}
    \label{fig:loss-steps-splice}
\end{figure*}
\section{Numerical Experiments}\label{4}
\label{sec:exps}
To determine whether the proposed variant yields practical advantages, we consider 5 different convex optimization problems of the form \eqref{eq:main_problem}, incorporating both synthetic and real-world data. Specifically, we address
\begin{equation}
    \min_{\bx\in\R^\m} f(\bx) + \varphi(\bx) := \frac{1}{N}\sum_{j=1}^N \ell(\bx, z_j) + \varphi(\bx)
\end{equation}
where $z_j = (a_j, b_j)\in\R^\m\times\R$ represents an input-output pair drawn from an (unknown) underlying distribution $\mathcal{D}$ and $N$ is the number of samples.
Following the above considerations, we employ datasets of the form $(A, \boldsymbol{b})$ with $A \in \R^{N \times \m}$ and $\boldsymbol{b} \in \R^{N}$. We compare our algorithm with AdaGrad over 10 independent initializations. We use for both methods a small constant $\varepsilon = 10^{-8}$, for numerical stability.  
% \textbf{Robustness with respect to the stepsize parameter $\eta$.}
%
% \MB{We study the influence of the stepsize parameter $\eta$ on the performance of both methods. The results reported in the first row of Figures \ref{fig:robustness_nonsmooth} and \ref{fig:robustness_smooth} show that AdaGrad-Diff converges reliably across a much wider range of learning rates compared to standard AdaGrad. In particular, for moderate choices of $\eta$, the method consistently achieves faster convergence. Even in extreme regimes, when $\eta$ is either very small or very large, AdaGrad-Diff maintains meaningful progress along the optimization trajectory, reaching points closer to the optimum than its counterpart. This demonstrates that the proposed accumulation of gradient differences stabilizes the dynamics and enhances robustness to hyperparameter selection, reducing the need for careful tuning.}
% Figure \ref{fig:loss-steps-splice} further shows that, when $\eta$ is close to the optimally tuned ones, the two methods exhibit comparable performance. However, for suboptimal choices of $\eta$, AdaGrad-Diff consistently yields significantly improved performance compared to AdaGrad. These observations were consistently reproduced across all experiments presented later in this section, suggesting a robust empirical trend. Due to space constraints, these results are provided in Appendix \ref{B}.
% \textbf{Minimization performance.}
% We also evaluate the minimization performance of AdaGrad and AdaGrad-Diff using optimally tuned values of the stepsize parameter $\eta$ for the experiments considered in this section. As shown in the second row of Figures \ref{fig:robustness_nonsmooth} and \ref{fig:robustness_smooth}, AdaGrad-Diff achieves convergence behavior that is comparable to, and in several cases better than, that of AdaGrad.
In the following, both settings provided by Assumption \ref{ass:1} and Assumption \ref{ass:2} are considered.
\subsection{The $G$-Lipschitz continuous setting}
We study two problems satisfying Assumption~\ref{ass:1}, using either the \textbf{Hinge loss} or the \textbf{Least Absolute Deviation (LAD) loss} with $\ell_1$-regularization:
\begin{align*}
    & \ell(\bx, (a, b)) = 
    \begin{cases} 
    \max(0, 1 - b \langle \bx, a \rangle) & \text{(Hinge loss)}\\
    |b - \langle \bx, a \rangle| & \text{(LAD loss)}
    \end{cases} \\
    & \reg(\bx) = \lambda \lVert\bx\rVert_1
\end{align*}
where $a \in \R^\m$ and $b \in \R$ and $\lambda > 0$ is the regularization parameter. 
We generate synthetic datasets with all components drawn independently from a standard normal distribution. Specifically, the matrix $A$ has i.i.d.\ $\mathcal{N}(0,1)$ entries, as do the vector $w \in \R^\m$ and the noise $\epsilon\in \R^N$. To induce sparsity, a mask is applied to the entries of $w$. The response vector is, namely, either
\begin{equation*}
    \boldsymbol{b} = \text{sgn}(Aw + \epsilon)\quad \text{or}\quad\boldsymbol{b} = Aw + \epsilon ,
\end{equation*}
where $\text{sgn}(\cdot)$ denotes the sign function. We set $N = 500$, $\m = 100$, and $\lambda = 10^{-2}$, and the vector $w$ has 20 nonzero entries. A budget of 1000 iterations is considered for both problems.
\subsection{The $L$-Lipschitz smooth setting}
We consider three binary classification problems satisfying Assumption \ref{ass:2}. \\

\noindent
\textbf{Logistic Regression with $\ell_2$-penalty}. In this first problem, we incorporate the $\ell_2$ term into the smooth term $f$. Namely, we set 
\begin{equation*}
    \begin{cases}
        f(\bx) = \frac{1}{N}\sum_{j=1}^N\log(1 + \exp(-b_j\scalarp{a_j,\bx})) + \frac{\sigma}{2}\lVert\bx\rVert_2^2 \\
        \reg \equiv 0
    \end{cases}  
\end{equation*}
with $\sigma > 0$ the regularization parameter, $\bx \in \R^\m, b_j \in \{+1, -1\}$. 
The dataset used in this experiment is the \texttt{news20.binary}\footnote{ \url{https://www.csie.ntu.edu.tw/~cjlin/libsvmtools/datasets/} \label{libsvm}} dataset from the LIBSVM collection, which consists of $N = 15935$ samples in dimension $\m = 62061$. The regularization parameter is set to $\sigma = 10^{-4}$, and a budget of 2500 iterations is considered.\\

\noindent
\textbf{Logistic Regression with $\ell_1$-penalty}.
In this problem, we employ an $\ell_1$-regularization, so that the nonsmooth term is captured by $\reg$. Specifically, 
\begin{equation*}
\begin{cases}
    \ell(\bx, (a, b)) = \log(1 + \exp(-b\scalarp{a,\bx})) \\
    \reg(\bx) = \lambda \lVert \bx \rVert_1.
\end{cases}  
\end{equation*}
The dataset we consider now is the \texttt{splice.t} dataset\footref{libsvm} with $N = 2175$, $\m = 60$.
We set $\lambda = 10^{-2}$, and we consider a budget of 1000 iterations.\\

\noindent
\textbf{Support Vector Machines (SVM) for Classification.} We consider a well-known classification problem based on the hinge loss, Gaussian kernels, and $\ell_2$-regularization. Let $\Lambda: \R^\m \to \mathcal{H}$ denote the feature map associated with the Gaussian kernel, where $\mathcal{H}$ is the corresponding reproducing kernel Hilbert space (RKHS). The learning task consists of minimizing a regularized empirical risk functional in the primal formulation, namely
\begin{equation}
\min_{h\in \mathcal{H}} \;
\frac{1}{N} \sum_{j=1}^N \max\left(0, 1 -b_j\scalarp{h,\Lambda(a_j)}\right)  
+ \frac{\lambda}{2}\lVert h\rVert_\mathcal{H}^2
\end{equation}
We focus on the dual formulation of this problem. Denoting by $\balpha \in \mathbb{R}^N$ the vector of dual variables, the dual problem is defined by setting
\begin{equation*}
\begin{cases}
    f(\balpha) = \frac{1}{2\lambda}\balpha^\top K \balpha - \langle \balpha  , \mathbf{b} \rangle  \\
    \reg(\balpha) = \iota_{[0,1/N]^N}(\mathbf{b} \odot \balpha).
\end{cases}
\end{equation*}
Here, $K \in \R^{N \times N}$ denotes the kernel matrix, and $\odot$ is the Hadamard product. The indicator function $\iota_{[0,1/N]^N}$ enforces the box constraints $0 \le b_j\alpha_j \le 1/N$ for all $j \in [N]$.
We consider the 2-dimensional synthetic dataset \texttt{2moons}, which is available via the scikit learn library \cite{scikit}. We choose $N = 300$.
Moreover, we set $\lambda = 10^{-3}$ and the Gaussian width equal to 1. A budget of 20 iterations is considered.

\subsection{Effect of the stepsize parameter $\eta$}
\textbf{Robustness to $\eta$.}
%with respect to the stepsize parameter $\eta$.} 
We study the influence of the stepsize parameter $\eta$ on the performance of both methods. The results reported in the first row of Figures \ref{fig:robustness_nonsmooth} and \ref{fig:robustness_smooth} show that AdaGrad-Diff converges reliably across a much wider range of choices for $\eta$ compared to standard AdaGrad. In particular, for moderate choices of $\eta$, the method consistently achieves faster convergence. Even in extreme regimes, when $\eta$ is either very small or very large, AdaGrad-Diff maintains meaningful progress along the optimization trajectory, reaching lower objective values than its counterpart. This demonstrates that the proposed accumulation of gradient differences stabilizes the dynamics and enhances robustness to hyperparameter selection, reducing the need for careful tuning.\\

\noindent
Figure \ref{fig:loss-steps-splice} further shows that, when $\eta$ is close to the optimally tuned ones, the two methods exhibit comparable performance. However, for suboptimal choices of $\eta$, AdaGrad-Diff consistently yields significantly improved performance compared to AdaGrad. These observations were consistently reproduced across all experiments considered in the section, suggesting a robust empirical trend. Due to space constraints, these results are provided in Appendix \ref{B}.\\

\noindent
\textbf{Minimization performance.}
We also evaluate the minimization performance of AdaGrad and AdaGrad-Diff using optimally tuned values of the stepsize parameter $\eta$ for the experiments considered in this section. As shown in the second row of Figures \ref{fig:robustness_nonsmooth} and \ref{fig:robustness_smooth}, AdaGrad-Diff achieves convergence behavior that is comparable to, and in several cases better than, that of AdaGrad.

\section{Discussion}
We proposed an AdaGrad-style adaptive variant, where the adaptation is driven not by the cumulative (sub)gradient squared norms but by the cumulative squared norms of successive (sub)gradient differences. 
Empirically, we showed across several problems that the proposed method results in increased robustness with respect to the stepsize parameter $\eta$ when compared to AdaGrad, reducing sensitivity to hyperparameter tuning, a highly desirable property in practice. 

% We note that our derivation is of independent interest for advancing parameter-free optimization methods.
% In both \cite{ivgi2023dog} and \cite{defazio2023learning}, the authors aim to construct a parameter-free variant of AdaGrad in which the sequence $\eta_\iter$ serves as an online estimate of the distance between the initial point $\bx^1$ and a minimizer $\bx_\star$. However, this estimate is non-increasing, and consequently, the term $\sum_{\ite=1}^\iter \eta_\ite \bigl( \reg(\bx^{\ite+1}) - \reg(\bx^\ite) \bigr)$ obtained when summing over the iterates in \eqref{AdaGradBound}, is not necessarily non-positive, leaving its control unclear, unlike in Remark \ref{rmk:adagrad_bound}. To circumvent this issue, they omit the regularization term altogether and develop their theory in the unregularized setting.

\subsection{Limitations}
Even if standard in AdaGrad's analysis, a limitation of our result in the nonsmooth case is that the iterates are required to be bounded. This is not a problem in the context of a bounded constraint set, but in general, it is a condition that cannot be ensured a priori.
 A further technical aspect concerns the form of the final bound. Specifically, in the last inequality of the proofs of Theorems~\ref{thm_nonsmooth} and  \ref{thm_smooth}, a term appears that depends on the inverse of the initial gradient components. While this dependence may seem counterintuitive, it is not specific to our analysis and also arises in the Dual Averaging variant of AdaGrad \cite{duchi2011adaptive} as well as in more recent work \cite{defazio2023learning}. 

% In particular, since implementations typically choose a small parameter $\varepsilon \ll 1$ to avoid stalling at early iterations, the resulting bounds may involve very large constants when some entries of the initial gradient are extremely small or zero. While well-known in the literature, we regard this sensitivity as an unsatisfactory aspect of the current analysis.

% Our theoretical analysis is restricted to deterministic convex optimization. Extending the results to either nonconvex and/or stochastic settings remains an important open direction. Moreover, a technical limitation in the nonsmooth setting is the need to assume bounded iterates, an assumption that is standard in AdaGrad analyses and shared by our approach. Finally, a minor issue arises in the form of the final bound. Indeed, in the last inequality in the proof of Theorem \ref{thm_nonsmooth}, we obtain a term that depends on the inverse of the initial gradient components. This appears counterintuitive, although this issue is not new, as it also appears in the Dual Averaging variant of AdaGrad \cite{duchi2011adaptive} and in more recent analyses \cite{defazio2023learning}.  

\subsection{Future Work}
We conclude by outlining several directions that we consider both important and challenging.
%\textbf{Nonconvex optimization.}  
%Extending the analysis of AdaGrad-
First, extending the analysis of AdaGrad-Diff to nonconvex objectives is a natural and practically relevant, as many learning problems are inherently nonconvex. Second, while our results focus on deterministic optimization, extending them to stochastic settings under weaker and more realistic assumptions than almost-sure gradient boundedness remains an open challenge. Third, given the central role of Adam and related methods, exploring AdaGrad-Diff–style modifications incorporating exponential moving averages and momentum is a promising direction. Finally, investigating parameter-free variants may reveal additional benefits of difference-based stepsize mechanisms.


%\textbf{Stochastic optimization.} The results presented in this work focus on deterministic optimization. Existing approaches, such as those in \cite{carmon2022making,ivgi2023dog, defazio2023learning}, can be used to extend our nonsmooth convex analyses to stochastic settings with high-probability guarantees, but such arguments rely on almost-sure boundedness of the stochastic gradients.  Developing guarantees under weaker and more realistic assumptions remains an open problem.
%\textbf{Exponential moving averages and momentum.} Given the central role of Adam and its variants in modern machine learning, a promising direction is to explore AdaGrad-Diff–style modifications with exponential moving average and momentum.
%\textbf{Parameter-free optimization.} A final direction for future work is to investigate parameter-free strategies, in the spirit of \cite{ivgi2023dog} and \cite{defazio2023learning}, to assess whether the proposed difference-based stepsize mechanism can yield additional benefits in such settings.
\label{sec:disc}
%\subsection{Tables}
%You may also want to include tables that summarize material. Like figures, these should be centered, legible, and numbered consecutively. However, place the title \emph{above} the table with at least 0.1~inches of space before the title and the same after it, as in \cref{sample-table}. The table title should be set in 9~point type and centered unless it runs two or more lines, in which case it should be flush left.

%% Note use of \abovespace and \belowspace to get reasonable spacing
%% above and below tabular lines.

%\begin{table}[t]
%  \caption{Classification accuracies for naive Bayes and flexible Bayes on various data sets.}
%  \label{sample-table}  \begin{center}    \begin{small}      \begin{sc} \begin{tabular}{lcccr}          \toprule          Data set  & Naive         & Flexible      & Better?  \\        \midrule        Breast    & 95.9$\pm$ 0.2 & 96.7$\pm$ 0.2 & $\surd$  \\        Cleveland & 83.3$\pm$ 0.6 & 80.0$\pm$ 0.6 & $\times$ \\        Glass2    & 61.9$\pm$ 1.4 & 83.8$\pm$ 0.7 & $\surd$  \\Credit    & 74.8$\pm$ 0.5 & 78.3$\pm$ 0.6 &          \\       Horse     & 73.3$\pm$ 0.9 & 69.7$\pm$ 1.0 & $\times$ \\  Meta      & 67.1$\pm$ 0.6 & 76.5$\pm$ 0.5 & $\surd$  \\ Pima      & 75.1$\pm$ 0.6 & 73.9$\pm$ 0.5 &          \\     Vehicle   & 44.9$\pm$ 0.6 & 61.5$\pm$ 0.4 & $\surd$  \\   \bottomrule      \end{tabular}\end{sc}  \end{small} \end{center} \vskip -0.1in \end{table}

%Tables contain textual material, whereas figures contain graphical material. Specify the contents of each row and column in the table's topmost row. Again, you may float tables to a column's top or bottom, and set wide tables across both columns. Place two-column tables at the top or bottom of the page.

%\subsection{Theorems and Such}
%The preferred way is to number definitions, propositions, lemmas, etc. consecutively, within sections, as shown below. 
%\begin{definition}  \label{def:inj}  A function $f:X \to Y$ is injective if for any $x,y\in X$ different, $f(x)\ne    f(y)$. \end{definition}
%Using \cref{def:inj} we immediate get the following result:
%\begin{proposition} If $f$ is injective mapping a set $X$ to another set $Y$,  the cardinality of $Y$ is at least as large as that of $X$ \end{proposition}
%\begin{proof} Left as an exercise to the reader. \end{proof}
%\cref{lem:usefullemma} stated next will prove to be useful.
%\begin{lemma} \label{lem:usefullemma}  For any $f:X \to Y$ and $g:Y\to Z$ injective functions, $f \circ g$ is  injective. \end{lemma}
%\begin{theorem} \label{thm:bigtheorem}  If $f:X\to Y$ is bijective, the cardinality of $X$ and $Y$ are the same. \end{theorem} An easy corollary of \cref{thm:bigtheorem} is the following:
%\begin{corollary}  If $f:X\to Y$ is bijective, the cardinality of $X$ is at least as large as that of $Y$. \end{corollary}
%\begin{assumption} The set $X$ is finite. \label{ass:xfinite} \end{assumption}
%\begin{remark} According to some, it is only the finite case (cf. \cref{ass:xfinite}) that is interesting. \end{remark}
%restatable

%\subsection{Citations and References}
%Please use APA reference format regardless of your formatter or word processor. If you rely on the \LaTeX\/ bibliographic facility, use \texttt{natbib.sty} and \texttt{icml2026.bst} included in the style-file package to obtain this format.
%
%Citations within the text should include the authors' last names and year. If the authors' names are included in the sentence, place only the year in parentheses, for example when referencing Arthur Samuel's pioneering work \yrcite{Samuel59}. Otherwise place the entire reference in parentheses with the authors and year separated by a comma \cite{Samuel59}. List multiple references separated by semicolons \cite{kearns89,Samuel59,mitchell80}. Use the `et~al.' construct only for citations with three or more authors or after listing all authors to a publication in an earlier reference \cite{MachineLearningI}.

%Authors should cite their own work in the third person in the initial version of their paper submitted for blind review. Please refer to \cref{author info} for detailed instructions on how to cite your own papers.

%Use an unnumbered first-level section heading for the references, and use a hanging indent style, with the first line of the reference flush against the left margin and subsequent lines indented by 10 points. The references at the end of this document give examples for journal articles \cite{Samuel59}, conference publications \cite{langley00}, book chapters \cite{Newell81}, books \cite{DudaHart2nd}, edited volumes \cite{MachineLearningI}, technical reports \cite{mitchell80}, and dissertations \cite{kearns89}.

%Alphabetize references by the surnames of the first authors, with single author entries preceding multiple author entries. Order references for the same authors by year of publication, with the earliest first. Make sure that each reference includes all relevant information (e.g., page numbers).

%Please put some effort into making references complete, presentable, and consistent, e.g. use the actual current name of authors. If using bibtex, please protect capital letters of names and abbreviations in titles, for example, use \{B\}ayesian or \{L\}ipschitz in your .bib file.

%\section*{Accessibility}
%Authors are kindly asked to make their submissions as accessible as possible for everyone including people with disabilities and sensory or neurological differences. Tips of how to achieve this and what to pay attention to will be provided on the conference website \url{http://icml.cc/}.

%\section*{Software and Data} 
%If a paper is accepted, we strongly encourage the publication of software and data with the camera-ready version of the paper whenever appropriate. This can be done by including a URL in the camera-ready copy. However, \textbf{do not} include URLs that reveal your institution or identity in your submission for review. Instead, provide an anonymous URL or upload the material as ``Supplementary Material'' into the OpenReview reviewing system. Note that reviewers are not required to look at this material when writing their review.

% Acknowledgements should only appear in the accepted version.
%\section*{Acknowledgements}
%If a paper is accepted, the final camera-ready version can (and usually should) include acknowledgements.  Such acknowledgements should be placed at the end of the section, in an unnumbered section that does not count towards the paper page limit. Typically, this will include thanks to reviewers who gave useful comments, to colleagues who contributed to the ideas, and to funding agencies and corporate sponsors that provided financial support.

%\newpage 

% \section*{Impact Statement}
% This paper presents work whose goal is to advance the field of Machine
% Learning. There are many potential societal consequences of our work, none
% which we feel must be specifically highlighted here.

\bibliographystyle{plainnat}
\bibliography{main_icml}

%%%%%%%%%%%%%%%%%%%%%%%%%%%%%%%%%%%%%%%%%%%%%%%%%%%%%%%%%%%%%%%%%%%%%%%%%%%%%%%
%%%%%%%%%%%%%%%%%%%%%%%%%%%%%%%%%%%%%%%%%%%%%%%%%%%%%%%%%%%%%%%%%%%%%%%%%%%%%%%
% APPENDIX
%%%%%%%%%%%%%%%%%%%%%%%%%%%%%%%%%%%%%%%%%%%%%%%%%%%%%%%%%%%%%%%%%%%%%%%%%%%%%%%
%%%%%%%%%%%%%%%%%%%%%%%%%%%%%%%%%%%%%%%%%%%%%%%%%%%%%%%%%%%%%%%%%%%%%%%%%%%%%%%

\newtheorem{innercustomthm}{{\bf Theorem}}
\newenvironment{customthm}[1]
  {\renewcommand\theinnercustomthm{{\bf #1}}\innercustomthm}
  {\endinnercustomthm}

\newtheorem{innercustomlem}{{\bf Lemma}}
\newenvironment{customlem}[1]
  {\renewcommand\theinnercustomlem{{\bf #1}}\innercustomlem}
  {\endinnercustomlem}
  
  \newtheorem{innercustomcor}{{\bf Corollary}}
\newenvironment{customcor}[1]
  {\renewcommand\theinnercustomcor{{\bf #1}}\innercustomcor}
  {\endinnercustomcor}

\newtheorem{innercustomprop}{{\bf Proposition}}
\newenvironment{customprop}[1]
  {\renewcommand\theinnercustomprop{{\bf #1}}\innercustomprop}
  {\endinnercustomprop}


% \newpage
\appendix
% \onecolumn
\begin{center}
{\Large \bf Appendices} 
\end{center}

\section{Proofs of Section~\ref{sec:analysis}} 
\label{A}

In this section we set
\begin{equation*}
    (\forall\,\iter\in \N)(\forall\,i\in \{1,\dots, m\})\quad
    v_i^\iter:= \sqrt{\sum_{\ite=1}^\iter\norm{\uu_i^\ite- \uu_i^{\ite-1}}},
\end{equation*}
so that, $\w_i^\iter= \varepsilon + v_i^\iter$.

\begin{remark}
\label{rmk1}
Let $\HH$ be a Hilbert space and let $\W\colon \HH\to \HH$ be a positive self-adjoint operator which is bounded from below, meaning that $\W \succeq \alpha\Id$ for some $\alpha>0$. We define
\begin{equation*}
\scalarp{\,\cdot\,,\,\cdot\,}_\W\colon \HH\times\HH\to \R\colon (\x,\y)\mapsto \scalarp{\W\x,\y}.
\end{equation*}
We have that $\scalarp{\,\cdot\,,\,\cdot\,}_\W$ is a scalar product on $\HH$ and $(\HH, \scalarp{\,\cdot\,,\,\cdot\,}_\W)$ is also a Hilbert space. For the corresponding norm $\norm{\x}_\W = (\scalarp{\W\x,\x}_\W)^{1/2}$ we have
\begin{equation}
\label{eq:20260124a}
(\forall\,\x\in \HH)\quad \frac{1}{\norm{\W^{-1}}} \norm{\x}^2 \leq \norm{\x}_\W^2 \leq \norm{\W}\norm{\x}^2.
\end{equation}
Indeed $\norm{\x}_\W^2 = \scalarp{\W\x,\x} \leq \norm{\W} \norm{\x}^2$
and $\norm{\x}^2 = \scalarp{\W^{-1}\W \x, \x}\leq \norm{\W^{-1}} \norm{\x}_\W^2$.
Note that if $\HH = \bigoplus_{i=1}^m \HH_i$ and $\W = \bigoplus_{i=1}^m \w_i \Id_i$, with $w_i>0$ for every $i\in [m]$, then
\begin{equation*}
\norm{\W} = \sup_{\norm{\x}^2\leq 1} \scalarp{\W\x, \x} =
\sup \bigg\{ \sum_{i=1}^m \w_i \norm{\x_i}^2 \,\Big\vert\, \x\in \HH\ \text{and}\ \sum_{i=1}^m \norm{\x_i}^2\leq 1\bigg\} = \max_{1\leq i\leq m}{\w_i}.
\end{equation*}
Similarly $\norm{\W^{-1}} = \max_{1\leq i\leq m}\w_i^{-1} = (\min_{1\leq i\leq m} \w_i)^{-1}$.
Moreover, suppose that $f\in \Gamma_0(\HH)$ and let $\x\in \HH$. Then, for every $\uu\in \HH$,
\begin{equation*}
\uu\in \partial f(x)\ \Leftrightarrow\ (\forall\, \y\in \HH)\ f(\y) \geq f(\x)+\scalarp{\y-\x, \uu}
\ \Leftrightarrow\ (\forall\, \y\in \HH)\ f(\y) \geq f(\x)+\scalarp{\y-\x, \W^{-1}\uu}_{\W}.
\end{equation*}
Thus, we have $ \uu\in \partial f(\x)\ \Leftrightarrow\ \W^{-1}\uu\in \partial^{\W}\!f(\x)$, where $\partial^{\W}\!f(\x)$ denotes the subdifferential of $f$ at $\x$ in the Hilbert space $(\HH, \scalarp{\,\cdot\,, \,\cdot\,}_{\W})$. 
%Finally we note that if $f\colon \HH\to \R$ is Lipschitz smooth with constant $L$ in the Hilbert space $\HH$, then, using \eqref{eq:20260124a}, we have
%\begin{align*}
%(\forall\,\x,\y\in \HH)\quad
%\norm{\nabla^\W\! f(\x)- \nabla^\W\! f(\y)}_{\W}^2
%&= \norm{\W^{-1}\nabla f(\x)- \W^{-1}\nabla f(\y)}_{\W}^2 \\[1ex]
%&= \scalarp{\nabla f(\x)- \nabla f(\y), \W^{-1}(\nabla f(\x)- \nabla f(\y))}\\[1ex]
%&\leq \norm{\W^{-1}}\norm{\nabla f(\x)- \nabla f(\y)}^2\\[1ex]
%&\leq L^2\norm{\W^{-1}} \norm{\x-\y}^2\\[1ex]
%&\leq L^2\norm{\W^{-1}}^2 \norm{\x-\y}_{\W}^2.
%\end{align*}
%Thus, $\nabla^\W\! f$ is Lipschitz continuous with constant $L \norm{\W^{-1}}$ in the Hilbert space
%$(\HH, \scalarp{\,\cdot\,, \,\cdot\,}_{\W})$.
\end{remark}

\begin{lemma}
\label{lem:basicPGstep}
Let $X$ be a real Hilbert space,
let $\W\colon \HH\to \HH$ be a positive self-adjoint operator bounded from below and set
$\scalarp{\x,\y}_{\W} = \scalarp{\W\x,\y}$, for every $\x,\y\in \HH$.
Let $f,\reg\in \Gamma_0(X)$ and set $F = f+\reg$.
Let $x \in X$, $\uu \in \partial f(x)$, $\eta>0$, and set 
%$\x_+ = \argmin_{\y\in \HH} \big\{ \scalarp{\y-\x, \uu} + \reg(\y)+ (2\eta)^{-1}\norm{\y-\x}_{\W}^2 \big\}$.
$x_+ = \prox^\W_{\eta g}(x - \eta \W^{-1}\uu) = \argmin_{\y\in \HH} \big\{ \scalarp{\y-\x, \uu} + \reg(\y)+ (2\eta)^{-1}\norm{\y-\x}_{\W}^2 \big\}$. 
Then, for every $z\in \HH$
\begin{equation}
\label{eq:basicPGstep}
2\eta (F(x_+)- F(z)) \leq \norm{x-z}_{\W}^2 - \norm{x_+-z}_{\W}^2 - \norm{x_+-x}_{\W}^2 +2\eta \big[f(x_+)-f(x)- \pair{x_+-x}{\uu}\big].
\end{equation}
%In particular the following hold.
%\begin{enumerate}[label={\rm (\roman*)}]
%\item\label{lem:basicPGstep_i} Suppose that $\uu_+\in \partial f(x_+)$. Then
%\begin{equation*}
%(\forall\,z\in X)\quad 2\eta (F(x_+)- F(z)) \leq \norm{x-z}_{\W}^2 - \norm{x_+-z}_{\W}^2 +\eta^2 \norm{\uu - \uu_+}_{\W^{-1}}^2.
%\end{equation*}
%\item\label{lem:basicPGstep_ii} Suppose that $f\colon X\to \R$ is Lipschitz smooth with constant $L$. Then
%\begin{equation*}
%(\forall\,z\in X)\quad 2\eta (F(x_+)- F(z)) \leq \norm{x-z}_{\W}^2 - \norm{x_+-z}_{\W}^2 + \eta^2 \norm{\nabla f(x_+)-\nabla f(x)}^2
%- \frac{\eta}{L} \norm{\nabla f(x)-\nabla f(x_+)}^2.
%\end{equation*}
%\end{enumerate}
\end{lemma}
\begin{proof}
%As a preliminary observation we note that $(\HH, \scalarp{\,\cdot\,,\,\cdot\,}_\W)$ is also a Hilbert space.
%We first prove \eqref{eq:basicPGstep}. 
Let $z\in X$.
By definition of $x_+$ and Fermat's rule, $0\in \uu + \W(\x_+-\x)/\eta +  \partial \reg(\x_+)$
we have $\W(x-x_+)/\eta - \uu\in \partial \reg(x_+)$. Thus, by definition of subdifferential,
\begin{equation*}
\reg(z)-\reg(x_+) \geq \pair{z-x_+}{\W(x-x_+)/\eta - \uu} = \eta^{-1}\pair{z-x_+}{\W(x-x_+)}
-\pair{z-x_+}{\uu}.
\end{equation*}
Multiplying both side by $-2\eta$ we have
\begin{equation*}
2\eta(\reg(x_+)-\reg(z)) \leq  2\pair{x_+-z}{\W(x-x_+)}
+2\eta\pair{z-x_+}{\uu}.
\end{equation*}
Now we note that
\begin{equation*}
\norm{\x- z}_{\W}^2 = \norm{\x - x_+}_{\W}^2 + \norm{\x_+-z}_{\W}^2 +2\pair{\W(\x-x_+)}{\x_+-z}.
\end{equation*}
Hence
\begin{equation}
\label{eq:20250829a}
2\eta(\reg(\x_+)-\reg(z)) \leq  \norm{\x- z}_{\W}^2 - \norm{\x_+-z}_{\W}^2 -\norm{\x - \x_+}_{\W}^2 
+2\eta\pair{z-x_+}{\uu}.
\end{equation}
On the other hand, using the fact that $\uu \in \partial f(\x)$, it holds
\begin{equation*}
\pair{z-\x_+}{\uu} = \pair{z-\x}{\uu} + \pair{\x-\x_+}{\uu}
\leq f(z)-f(\x_+) + f(\x_+)- f(\x) - \pair{\x_+-\x}{\uu}.
\end{equation*}
Thus, 
\begin{equation}
\label{eq:20250829b}
2\eta \pair{z-\x_+}{\uu} \leq 2\eta (f(z)-f(\x_+)) + 2 \eta \big[f(\x_+)- f(\x) - \pair{\x_+-\x}{\uu}\big].
\end{equation}
Combining \eqref{eq:20250829a} and \eqref{eq:20250829b}, inequality \eqref{eq:basicPGstep} follows.
%
%\ref{lem:basicPGstep_i}: Referring to inequality \eqref{eq:basicPGstep}, we note that, being $\uu_+\in \partial f(x_+)$,
%\begin{equation*}
%f(x_+)-f(x)- \pair{x_+-x}{u} \leq \pair{x_+-x}{\uu_+}-\pair{x_+-x}{\uu} = \pair{x_+-x}{\uu_+-\uu}.
%\end{equation*}
%Thus, 
%\begin{equation*}
%(\forall\,z\in X)\quad 2\eta (F(x_+)- F(z)) \leq \norm{x-z}_{\W}^2 - \norm{x_+-z}_{\W}^2 - \norm{x-x_+}_{\W}^2 +2\eta \pair{x-x_+}{\uu - \uu_+}.
%\end{equation*}
%On the other hand using Young-Fenchel inequality we have
%\begin{align*}
%2\eta \pair{x-x_+}{\uu - \uu_+}-\norm{x-x_+}_{\W}^2
%&= 2\eta \pair{x-x_+}{\W^{-1}(\uu - \uu_+)}_{\W}-\norm{x-x_+}_{\W}^2\\[1ex]
%& =2 \bigg(\pair{x-x_+}{\eta\W^{-1}(\uu - \uu_+)}_{\W} -\frac{\norm{x-x_+}_{\W}^2}{2}\bigg)\\
%&\leq 2 \frac{\norm{\eta\W^{-1}(\uu -\uu_+)}_{\W}^2}{2} = \eta^2\norm{\uu- \uu_+}_{\W^{-1}}^2
%\end{align*}
%and the statement follows.
%
%\ref{lem:basicPGstep_ii}: 
%By the Lipschitz smoothness of $f$ we have
%\begin{equation*}
%\frac{1}{2L}\norm{\nabla f(x_+)-f(x)}^2\leq f(x)-f(x_+)- \pair{x-x_+}{\nabla f(x_+)}. 
%\end{equation*}
%Thus, from \eqref{eq:basicPGstep} (with $\uu= \nabla f(\x)$) we have
%\begin{align*}
%2\eta (F(x_+)- F(z)) &\leq \norm{x-z}_{\W}^2 - \norm{x_+-z}_{\W}^2 - \norm{x_+-x}_{\W}^2 +2\eta \big[f(x_+)-f(x)\big] - 2\eta\pair{x_+-x}{\nabla f(x)}\\[1ex]
%&\leq\norm{x-z}_{\W}^2 - \norm{x_+-z}_{\W}^2 - \norm{x_+-x}_{\W}^2\\
%&\qquad +2\eta \Big[\pair{x_+-x}{\nabla f(x_+)}- \frac{1}{2L} \norm{\nabla f(x)-\nabla f(x_+)}^2\Big] - 2\eta\pair{x_+-x}{\nabla f(x)}\\[1ex]
%& =\norm{x-z}_{\W}^2 - \norm{x_+-z}_{\W}^2\\
%&\hspace{15ex} + 2\Big[\pair{x_+-x}{\eta\W^{-1}(\nabla f(x_+)-\nabla f(x))}_{\W} - \frac1 2 \norm{x_+-x}_{\W}^2\Big]
%- \frac{\eta}{L} \norm{\nabla f(x)-\nabla f(x_+)}^2\\[1ex]
%&\leq \norm{x-z}_{\W}^2 - \norm{x_+-z}_{\W}^2 +  \norm{\eta\W^{-1}(\nabla f(x_+)-\nabla f(x))}_{\W}^2
%- \frac{\eta}{L} \norm{\nabla f(x)-\nabla f(x_+)}^2\\[1ex]
%&\leq \norm{x-z}_{\W}^2 - \norm{x_+-z}_{\W}^2 +  \eta^2\norm{\nabla f(x_+)-\nabla f(x)}_{\W^{-1}}^2
%- \frac{\eta}{L} \norm{\nabla f(x)-\nabla f(x_+)}^2.
%\end{align*}
%and the statement follows.
\end{proof}

%\begin{remark}
%\label{rmk2}
%Recalling Remark~\ref{rmk1}, let $f,\reg \in \Gamma_0(\HH)$, $\x\in \HH$, $\uu\in \partial f(\x)$, and set
%\begin{equation*}
%\x_+ = \prox^{\W}_{\eta \reg} (\x- \eta \W^{-1}\uu) 
%= \argmin_{\y\in \HH}\Big\{ \scalarp{\y-\x, \uu} + \reg(\y) + \frac{1}{2\eta} \norm{\y-\x}^2_{\W}\Big\}.
%\end{equation*}
%Applying Lemma~\ref{lem:basicPGstep} in the Hilbert space $(\HH, \scalarp{\,\cdot\,, \,\cdot\,}_{\W})$ the following hold
%\begin{enumerate}[label={\rm (\roman*)}]
%\item\label{rmk2_i} If $\uu_+\in \partial f(x_+)$, then, for every $z\in \HH$
%\begin{align}
%\label{eq:rmk2_i}
%\nonumber 2\eta (F(x_+)- F(z)) &\leq \norm{x-z}_\W^2 - \norm{x_+-z}_\W^2 +\eta^2 \norm{\W^{-1}(\uu - \uu_+)}_W^2\\[1ex]
%&= \norm{x-z}_\W^2 - \norm{x_+-z}_\W^2 +\eta^2 \norm{\uu - \uu_+}_{\W^{-1}}^2.
%\end{align}
%\item\label{lem:basicPGstep_ii} If $f\colon X\to \R$ is Lipschitz smooth with constant $L$ (in the Hilbert space $\HH$), then for every $z\in \HH$,
%\begin{align}
%\label{eq:rmk2_ii}
%\nonumber 2\eta (F(x_+)- F(z)) &\leq \norm{x-z}_\W^2 - \norm{x_+-z}_\W^2 + \eta^2 \norm{\W^{-1}(\nabla f(x_+)-\nabla f(x))}_\W^2\\[1ex]
%\nonumber&\hspace{15ex} - \frac{\eta}{L\norm{\W^{-1}}} \norm{\W^{-1}(\nabla f(x)-\nabla f(x_+))}_{\W}^2\\
%&\leq \norm{x-z}_\W^2 - \norm{x_+-z}_\W^2 + \eta^2 \norm{\nabla f(x_+)-\nabla f(x)}_{\W^{-1}}^2 - \frac{\eta}{L} \norm{\nabla f(x)-\nabla f(x_+)}^2.
%\end{align}
%\end{enumerate}
%\end{remark}
\allowdisplaybreaks

\begin{customlem}{\ref{lemma1}}[\textbf{The basic inequalities}]
Let $\bx\in \bHH$ and $\iter \in \N$. Then
\begin{equation}%\label{lemma1_bound_app}
\tag{\ref{lemma1_bound}}    2\eta_\iter(F(\bx^{\iter+1}) - F(\bx))  
    \le \norm{\bx^\iter - \bx}_{\W_\iter}^2- \norm{\bx^{\iter+1} - \bx}_{\W_\iter}^2+ \eta_\iter^2 \norm{\buu^{\iter+1} - \buu^{\iter}}_{\W_\iter^{-1}}^2.
    \end{equation}
Moreover, under Assumption~\ref{ass:2}, it holds
\begin{equation}%\label{lemma1_bound2_app}
\tag{\ref{lemma1_bound2}}    2\eta_\iter(F(\bx^{\iter+1}) - F(\bx))  
    \le \norm{\bx^\iter - \bx}_{\W_\iter}^2- \norm{\bx^{\iter+1} - \bx}_{\W_\iter}^2 + \eta_\iter^2 \norm{\buu^{\iter+1} - \buu^{\iter}}_{\W_\iter^{-1}}^2 + \frac{\eta_\iter}{L}\norm{\buu^{\iter+1} - \buu^{\iter}}^2.
    \end{equation}
\end{customlem}
\begin{proof}
We note that
\begin{equation*}
\bx^{\iter+1} = \prox^{\W_\iter}_{\eta_\iter \reg} (\bx_\iter- \eta_\iter \W_\iter^{-1}\buu_\iter) 
= \argmin_{\bx\in \bHH}\Big\{ \scalarp{\bx-\bx^\iter_\iter, \buu_\iter} + \reg(\bx) + \frac{1}{2\eta_\iter} \norm{\bx-\bx_\iter}^2_{\W_\iter}\Big\}.
\end{equation*}
Thus, 
%recalling Remark~\ref{rmk1}, and
applying Lemma~\ref{lem:basicPGstep} with $\W= \W_\iter$,
%in the Hilbert space $(\HH, \scalarp{\,\cdot\,, \,\cdot\,}_{\W_\iter})$, 
we have, for every $\bx\in \bHH$,
\begin{align}
\label{eq:20260124d}
    \nonumber 2\eta_\iter (F(\bx^{\iter+1})- F(\bx)) & \leq \norm{\bx^\iter\!-\bx}_{\W_\iter}^2\! - \norm{\bx^{\iter+1}-\bx}_{\W_\iter}^2 \!- \norm{\bx^{\iter+1}\!-\bx}_{\W_\iter}^2 \\ & \qquad +2\eta_\iter \big[f(\bx^{\iter+1})-f(\bx^\iter)- \pair{\bx^{\iter+1}\!-\bx^\iter}{\buu^\iter}\big].
\end{align}
Moreover, the following hold. 
\begin{enumerate}[label={\rm (\roman*)}]
\item\label{rmk2_i} Referring to inequality \eqref{eq:20260124d}, we note that, being $\buu^{\iter+1}\in \partial f(\bx^{\iter+1})$,
\begin{align*}
f(\bx^{\iter+1})-f(\bx^\iter)- \pair{\bx^{\iter+1}-\bx^\iter}{\buu^n} & \leq \pair{\bx^{\iter+1}-\bx^\iter}{\uu^{\iter+1}}-\pair{\bx^{\iter+1}-\bx^\iter}{\buu^{\iter}} \\
&= \pair{\bx^{\iter+1}-\bx^\iter}{\buu^{\iter+1}-\buu^\iter}.
\end{align*}
Thus, for every $\bx\in \bHH$
\begin{align*}
2\eta_\iter (F(\bx^{\iter+1})- F(\bx)) & \leq \norm{\bx^\iter-\bx}_{\W_\iter}^2 - \norm{\bx^{\iter+1}-\bx}_{\W_\iter}^2 - \norm{\bx^\iter-\bx^{\iter+1}}_{\W_\iter}^2 \\
& \qquad +2\eta_\iter 
\pair{\bx^\iter-\bx^{\iter+1}}{\buu^\iter - \buu^{\iter+1}}.
\end{align*}
On the other hand using Young-Fenchel inequality we have
\begin{align*}
2\eta_{\iter} \pair{\bx^\iter-\bx^{\iter+1}}{\buu^\iter - \buu^{\iter+1}}&-\norm{\bx^\iter-\bx^{\iter+1}}_{\W_\iter}^2 \\[1ex]
& = 2\eta_\iter \pair{\bx^\iter-\bx^{\iter+1}}{\W_\iter^{-1}(\buu^\iter - \buu^{\iter+1})}_{\W_\iter}-\norm{\bx^\iter-\bx^{\iter+1}}_{\W_\iter}^2\\[1ex]
& =2 \bigg(\pair{\bx^\iter-\bx^{\iter+1}}{\eta_\iter\W_{\iter}^{-1}(\buu^\iter - \buu^{\iter+1})}_{\W_\iter} 
-\frac{\norm{\bx^\iter-\bx^{\iter+1}}_{\W_\iter}^2}{2}\bigg)\\
&\leq 2 \frac{\norm{\eta_\iter\W_\iter^{-1}(\buu^\iter -\buu^{\iter+1})}_{\W_\iter}^2}{2} = \eta_\iter^2\norm{\buu^\iter- \buu^{\iter+1}}_{\W_\iter^{-1}}^2
\end{align*}
and \eqref{lemma1_bound} follows.
%Since $\uu_{\iter+1}\in \partial f(\x_\iter)$, then, for every $\x\in \HH$
%\begin{align}
%\label{eq:rmk2_i}
%\nonumber 2\eta_{\iter} (F(x^{\iter+1})- F(\x)) &\leq \norm{\x^\iter-\x}_{\W_\iter}^2 - \norm{\x^{\iter+1}-\x}_{\W_\iter}^2 +\eta_{\iter}^2 \norm{\W_\iter^{-1}(\uu^{\iter} - \uu^{\iter+1})}_{\W_\iter}^2\\[1ex]
%&= \norm{\x^\iter-\x}_{\W_\iter}^2 - \norm{\x^{\iter+1}-\x}_{\W_\iter}^2 +\eta_{\iter}^2 \norm{\uu^{\iter} - \uu^{\iter+1}}_{\W_\iter^{-1}}^2.
%\end{align}
\item\label{lem:basicPGstep_ii} 
Concerning the second part of the statement, suppose that $f:\bHH\to \R$ is Lipschitz smooth with constant $L$ (in the Hilbert space $\bHH$). 
The Lipschitz smoothness of $f$ implies that
\begin{equation*}
\frac{1}{2L}\norm{\nabla f(\bx^{\iter+1})-f(\bx^\iter)}^2\leq f(\bx^\iter)-f(\bx^{\iter+1})- \pair{\bx^\iter-\bx^{\iter+1}}{\nabla f(\bx^{\iter+1})}. 
\end{equation*}
Thus, from \eqref{eq:20260124d} (with $\buu^\iter= \nabla f(\bx^\iter)$) we have
\begin{align*}
2\eta_\iter (F(\bx^{\iter+1})- F(\bx)) &\leq \norm{\bx^\iter-\bx}_{\W_\iter}^2 - \norm{\bx^{\iter+1}-\bx}_{\W_\iter}^2 - \norm{\bx^{\iter+1}-\bx^{\iter}}_{\W_\iter}^2\\[1ex]
&\qquad +2\eta_\iter \big[f(\bx^{\iter+1})-f(\bx^{\iter})\big] - 2\eta_\iter\pair{\bx^{\iter+1}-\bx^{\iter}}{\nabla f(\bx^{\iter})}\\[1ex]
&\leq\norm{\bx^{\iter}-\bx}_{\W_\iter}^2 - \norm{\bx^{\iter+1}-\bx}_{\W_\iter}^2 - \norm{\bx^{\iter+1}-\bx^{\iter}}_{\W_\iter}^2\\
&\qquad +2\eta_\iter \Big[\pair{\bx^{\iter+1}-\bx^{\iter}}{\nabla f(\bx^{\iter+1})}- \frac{1}{2L} \norm{\nabla f(\bx^{\iter})-\nabla f(\bx^{\iter+1})}^2\Big] \\
&\qquad- 2\eta_\iter\pair{\bx^{\iter+1}-\bx^{\iter}}{\nabla f(\bx^{\iter})}\\[1ex]
& =\norm{\bx^{\iter}-\bx}_{\W_\iter}^2 - \norm{\bx^{\iter+1}-\bx}_{\W_\iter}^2\\
&\qquad + 2\pair{\bx^{\iter+1}-\bx^{\iter}}{\eta_\iter\W_\iter^{-1}(\nabla f(\bx^{\iter+1})-\nabla f(\bx^{\iter}))}_{\W_\iter} - \norm{\bx^{\iter+1}-\bx^{\iter}}_{\W_\iter}^2\\
&\qquad
- \frac{\eta_\iter}{L} \norm{\nabla f(\bx^{\iter})-\nabla f(\bx^{\iter+1})}^2\\[1ex]
&\leq \norm{\bx^{\iter}-\bx}_{\W_\iter}^2 - \norm{\bx^{\iter+1}-\bx}_{\W_\iter}^2 +  \norm{\eta_\iter\W_\iter^{-1}(\nabla f(\bx^{\iter+1})-\nabla f(\bx^{\iter}))}_{\W_\iter}^2\\[1ex]
&\qquad- \frac{\eta_\iter}{L} \norm{\nabla f(\bx^{\iter})-\nabla f(\bx^{\iter+1})}^2\\[1ex]
&\leq \norm{\bx^{\iter}-\bx}_{\W_\iter}^2 - \norm{\bx^{\iter+1}-\bx}_{\W_\iter}^2 +  \eta_\iter^2\norm{\nabla f(\bx^{\iter+1})-\nabla f(\bx^{\iter})}_{\W_\iter^{-1}}^2
\\
& \qquad -\frac{\eta_\iter}{L} \norm{\nabla f(\bx^{\iter})-\nabla f(\bx^{\iter+1})}^2,
\end{align*}
where in the penultimate equation we used Young-Fenchel inequality.
\end{enumerate}
%Then for every $\x\in \HH$,
%\begin{align}
%\label{eq:rmk2_ii}
%\nonumber 2\eta_\iter (F(\x^{\iter+1})- F(\x)) &\leq \norm{\x^{\iter}-\x}_{\W_\iter}^2 - \norm{\x^{\iter+1}-\x}_{\W_\iter}^2 + \eta_\iter^2 \norm{\W_\iter^{-1}(\nabla f(\x^{\iter+1})-\nabla f(\x^\iter))}_{\W_\iter}^2\\[1ex]
%\nonumber&\hspace{15ex} - \frac{\eta_\iter}{L\norm{\W_\iter^{-1}}} \norm{\W_\iter^{-1}(\nabla f(\x^\iter)-\nabla f(\x^{\iter+1}))}_{\W}^2\\
%\nonumber&\leq \norm{\x^\iter-\x}_{\W_\iter}^2 - \norm{\x^{\iter+1}-\x}_{\W_\iter}^2 + \eta_\iter^2 \norm{\nabla f(\x^{\iter+1})-\nabla f(\x^\iter)}_{\W_\iter^{-1}}^2\\
%\nonumber&\hspace{15ex} - \frac{\eta_\iter}{L\norm{\W_\iter^{-1}}} \norm{\nabla f(\x^\iter)-\nabla f(\x^{\iter+1})}_{\W^{-1}}^2\\
%\nonumber&\leq \norm{\x^\iter-\x}_{\W_\iter}^2 - \norm{\x^{\iter+1}-\x}_{\W_\iter}^2 + \eta_\iter^2 \norm{\nabla f(\x^{\iter+1})-\nabla f(\x^\iter)}_{\W_\iter^{-1}}^2\\
%&\hspace{15ex} - \frac{\eta_\iter}{L} \norm{\nabla f(\x^\iter)-\nabla f(\x^{\iter+1})}^2,
%\end{align}
%\end{enumerate}
%where in the last inequality we used that, for every $z\in \HH$, $\norm{z}^2_{\W^{-1}} = \scalarp{\W^{-1}z,z}  \geq (\min_{1\leq i\leq m} \w_i^{-1}) \norm{\x}^2$ and that $\norm{\W^{-1}} = \max_{1\leq i\leq m} \w_i^{-1} = (\min_{1\leq i\leq m} \w_i)^{-1}$.
The statement follows.
\end{proof}

\begin{customcor}{\ref{Proposition1}}
Let $\bx\in \bHH$ and $\iter \in \N$ and set
%Under the same assumptions of Lemma~\ref{lemma1}, set 
$\Delta^{\iter-1}_{\max} = \max_{1\leq \ite\leq \iter-1} \max_{1\leq i\leq m} \norm{\x^\ite_i - x_i}^2$. Then
\begin{equation}%\label{Proposition1_bound_app}
\tag{\ref{Proposition1_bound}}
\sum_{\ite=2}^\iter\eta_{\ite-1}(F(\bx^\ite) - F(\bx)) \le \frac{\varepsilon}{2}\lVert \bx^1 - \bx\rVert^2 + \frac{1}{2}\Delta_{\max}^{\iter-1} \sum_{i=1}^m \w_i^{\iter-1} + \frac{1}{2}\sum_{\ite=2}^\iter\eta_{\ite-1}^2 \lVert \buu^{\ite} - \buu^{\ite-1} \rVert_{\W_{\ite-1}^{-1}}^2. 
%\label{sum_lemma2}
\end{equation}
\end{customcor}
\begin{proof}
Let $n\in \N$, with $n\geq 2$. It follows from Lemma~\ref{lemma1} that
\begin{equation}
\label{eq:20260122a}
    2\sum_{\ite=1}^{\iter-1} \eta_\ite (F(\bx^{\ite+1})- F(\bx))
    \leq \sum_{\ite=1}^{\iter-1} \norm{\bx^\ite-\bx}_{\W_\ite}^2 - \norm{\bx^{\ite+1}-\bx}_{\W_\ite}^2 + \sum_{\ite=1}^{\iter-1} \eta_\ite^2\norm{\buu^{\ite+1}-\buu^\ite}_{\W_\ite^{-1}}^2.
\end{equation}
The first term on the right-hand side can be bounded as follows
\begin{align*}
    \sum_{\ite=1}^{\iter-1} \norm{\bx^\ite-\bx}_{\W_\ite}^2 - \norm{\bx^{\ite+1}-\bx}_{\W_\ite}^2 
    &= \norm{\bx^1-\bx}_{\W_1}^2 + \sum_{\ite=2}^{\iter-1} \norm{\bx^\ite-\bx}_{\W_\ite}^2 - \sum_{\ite=1}^{\iter-2} \norm{\bx^{\ite+1}-\bx}_{\W_\ite}^2\\
    &\qquad-\norm{\bx^{\iter}-\bx}_{\W_{\iter-1}}^2\\
    &\leq \norm{\bx^1-\bx}_{\W_1}^2 + \sum_{\ite=1}^{\iter-2} \norm{\bx^{\ite+1}-\bx}_{\W_{\ite+1}}^2 - \norm{\bx^{\ite+1}-\bx}_{\W_\ite}^2\\
    &= \norm{\bx^1-\bx}_{\W_1}^2 + \sum_{\ite=1}^{\iter-2} \norm{\bx^{\ite+1}-\bx}_{\W_{\ite+1}-\W_\ite}^2\\
    &= \varepsilon \norm{\bx^1-\bx}^2+ \sum_{i=1}^m v_i^1\norm{\x_i^1 -\x_i}^2 \\
    &\qquad + \sum_{\ite=1}^{\iter-2}\sum_{i=1}^m(\w_i^{\ite+1}-\w_i^\ite)\norm{\x_i^{\ite+1}-\x_i}^2\\
    &\leq \varepsilon \norm{\bx^1-\bx}^2+ \Delta_{\max}^{\iter-1}\sum_{i=1}^m v_i^1 + \Delta_{\max}^{\iter-1}\sum_{i=1}^m\sum_{\ite=1}^{\iter-2}(v_i^{\ite+1}-v_i^\ite)\\
    &= \varepsilon \norm{\bx^1-\bx}^2+ \Delta_{\max}^{\iter-1}\sum_{i=1}^m v_i^{\iter-1}
\end{align*}
Combining the latter inequality and \eqref{eq:20260122a} the statement follows.
\end{proof}
\noindent
In order to provide the convergence rate, \citet{duchi2011adaptive} makes use of the following result.
\begin{lemma}\label{lemma2}
Let $(a_\iter)_{\iter\in \N}\in \R_+^\N$ and $\varepsilon > 0$. Then,
\begin{equation*}
    \sum_{\ite=1}^{\iter} \frac{a_\ite}{\varepsilon +\sqrt{ \sum_{j=1}^{\ite} a_j}} \le 2\sqrt{\sum_{\ite=1}^{\iter} a_\ite}.
\end{equation*}
\end{lemma}

\begin{customprop}{\ref{main_theorem}}
%\begin{proposition}\label{main_theorem}
%Under Assumption~\ref{ass:0}, 
%let $(\x^\iter)_{\iter\in \N}$ be the iterates generated Algorithm~\ref{adadiff}, 
%Let $\x\in \HH$. Then for every $\iter \in \N$
Let $(\bx^\iter)_{\iter\in \N}$ be generated by Algorithm~\ref{adadiff}. Let $\bx\in \bHH$, $\iter \in \N$ and set
$\Delta^{\iter-1}_{\max} = \max_{1\leq \ite\leq \iter-1} \max_{1\leq i\leq \m} \norm{\x^\ite_i - x_i}^2$.
Then
\begin{multline}
\tag{\ref{extra}}
 F(\bar{\bx}^{\iter}) - F(\bx)  \le \frac{\varepsilon}{2\eta(\iter-1)}\lVert \bx^1 - \bx\rVert^2
+\frac{1}{\iter-1}\left[\frac{1}{2\eta}\Delta^{\iter-1}_{\max} + \eta \right]\sum_{i=1}^m \sqrt{\sum_{\ite=1}^{\iter} \norm{\uu_i^\ite - \uu_{i}^{\ite-1}}^2} \\
+ \frac{\eta}{2(\iter-1)} \sum_{i=1}^m\sum_{\ite=2}^{\iter} \left(\frac{1}{\w_i^{\ite-1}} - \frac{1}{\w_i^\ite}\right)\norm{\uu_i^\ite - \uu_i^{\ite-1}}^2.
\end{multline}
\end{customprop}
\begin{proof}
We assumed that $\eta_\iter\equiv \eta$.
It follows from Corollary~\ref{Proposition1} and Lemma~\ref{lemma2} that
\begin{align*}
\sum_{\ite=2}^\iter(F(\bx^\ite) - F(\bx)) 
    &\leq \frac{\varepsilon}{2\eta}\lVert \bx^1 - \bx\rVert^2 + \frac{1}{2\eta}\Delta_{\max}^{\iter-1} \sum_{i=1}^m v_i^{\iter-1} + \frac{\eta}{2}\sum_{\ite=2}^\iter \lVert \buu^{\ite} - \buu^{\ite-1} \rVert_{\W_{\ite-1}^{-1}}^2 \\
    &= \frac{\varepsilon}{2\eta}\lVert \bx^1 - \bx\rVert^2 + \frac{1}{2\eta}\Delta_{\max}^{\iter-1} \sum_{i=1}^m v_i^{\iter-1} + \frac{\eta}{2}\sum_{\ite=2}^\iter \lVert \buu^{\ite} - \buu^{\ite-1} \rVert_{\W_{\ite}^{-1}}^2\\
    &\hspace{15ex}+ \frac{\eta}{2}\sum_{\ite=2}^\iter \lVert \buu^{\ite} - \buu^{\ite-1} \rVert_{\W_{\ite-1}^{-1}- \W_{\ite}^{-1}}^2\\
    &= \frac{\varepsilon}{2\eta}\lVert \bx^1 - \bx\rVert^2 + \frac{1}{2\eta}\Delta_{\max}^{\iter-1} \sum_{i=1}^m v_i^{\iter-1} + \frac{\eta}{2}\sum_{\ite=2}^\iter \sum_{i=1}^m \frac{1}{\varepsilon+ v_i^\ite}\lVert \uu_i^{\ite} - \uu_i^{\ite-1} \rVert^2\\
    &\hspace{15ex}+ \frac{\eta}{2}\sum_{\ite=2}^\iter \sum_{i=1}^m \bigg(\frac{1}{\varepsilon+ v_i^{\ite-1}}-\frac{1}{\varepsilon + v_i^\ite}\bigg)\norm{\uu_i^{\ite} - \uu_i^{\ite-1}}^2\\
    &\leq \frac{\varepsilon}{2\eta}\lVert \bx^1 - \bx\rVert^2 + \frac{1}{2\eta}\Delta_{\max}^{\iter-1} \sum_{i=1}^m \sqrt{\sum_{k=1}^{\iter-1}\norm{\uu_i^\ite-\uu_i^{\ite-1}}^2} + \eta \sum_{i=1}^m \sqrt{\sum_{\ite=1}^\iter \lVert \uu_i^{\ite} - \uu_i^{\ite-1}\rVert^2}\\
    &\hspace{15ex}+ \frac{\eta}{2}\sum_{\ite=2}^\iter \sum_{i=1}^m \bigg(\frac{1}{\varepsilon+ v_i^{\ite-1}}-\frac{1}{\varepsilon + v_i^\ite}\bigg)\norm{\uu_i^{\ite} - \uu_i^{\ite-1}}^2
    \end{align*}
    where in the last inequality we used Lemma~\ref{lemma2} with $a_k = \norm{\uu_i^\ite- \uu_i^{\ite-1}}^2$. Defining $\bar{\bx}^\iter = (\iter-1)^{-1}\sum_{\ite=2}^\iter \bx^\ite$, Jensen inequality yields
    \begin{equation*}
    F(\bar{\bx}^\iter)- F(\bx) \leq \frac{1}{\iter-1} \sum_{\ite=2}^\iter F(\bx^\ite)-F(\bx)
    \end{equation*}
    and the statement follows.
\end{proof}


\subsection{Proofs of subsection~\ref{sec:analysis_b}}

\begin{customprop}{\ref{proposition5}}
 We have that $\sum_{\iter=1}^\infty \lVert \buu^{\iter+1} - \buu^\iter \rVert^2<+\infty$.
\end{customprop}
\begin{proof}
Let 
\begin{equation}
I = \bigg\{ i\in [m]\,\Big\vert\, \exists\,\iter\in \N\ \text{s.t.}\ \frac{\eta}{\varepsilon+ v^\iter_i} \leq \frac{1}{2L}\bigg\}.
\end{equation}
Then, for every $i\in [m]\setminus I$ and every $n\in \N$, $\eta/(\varepsilon+ v^\iter_i)>1/(2L)$ and hence
\begin{equation*}
\sum_{\ite=1}^{\iter} \norm{\uu_i^{\ite+1} - \uu^{\ite}_i}^2 \leq 
\sum_{\ite=2}^{\iter+1} \norm{\uu_i^{\ite} - \uu^{\ite-1}_i}^2 \leq (v_i^{\iter+1})^2 < (2\eta L- \varepsilon)^2.
\end{equation*}
Therefore,
\begin{equation}
\label{eq:20260124b}
(\forall\,\iter\in \N)\quad
\sum_{\ite=1}^\iter \sum_{i\in [m]\setminus I} \norm{\uu_i^{\ite+1} - \uu^{\ite}_i}^2 =
\sum_{i\in [m]\setminus I}\sum_{\ite=1}^\iter \norm{\uu_i^{\ite+1} - \uu^{\ite}_i}^2 < m(2\eta L- \varepsilon)^2.
\end{equation}
Thus, it follows from \eqref{eq:20260124b} that if $I= \varnothing$,  
\begin{equation*}
(\forall\,\iter\in \N)\quad
\sum_{\ite=1}^\iter \norm{\buu^{\ite} - \buu^{\ite-1}}^2 = \sum_{\ite=1}^\iter \sum_{i\in [m]}
\norm{\uu_i^{\ite} - \uu^{\ite-1}_i}^2 \leq m(2\eta L- \varepsilon)^2
\end{equation*}
and the statement follows.
On the other hand, suppose that $I\neq\varnothing$ and set
\begin{equation*}
(\forall\,i\in I)\quad \iter_i = \min\bigg\{ \iter\in \N\,\Big\vert\, \frac{\eta}{\varepsilon+ v^\iter_i} \leq 
\frac{1}{2L}\bigg\}\quad\text{and}\quad \bar{\iter} = \max_{i\in I} \iter_i.
\end{equation*}
Let $\ite \in \N$ with $\ite\geq \bar{\iter}$.
By inequality \eqref{lemma1_bound2} in Lemma~\ref{lemma1} with $\iter=\ite$ and $\bx= \bx^{\ite}$ we have
\begin{align*}
    2(F(\bx^{\ite+1}) - F(\bx^\ite))  
    &\leq \eta \norm{\buu^{\ite+1} - \buu^{\ite}}_{\W_\ite^{-1}}^2 - \frac{1}{L}\norm{\buu^{\ite+1} - \buu^{\ite}}^2\\
    &=\sum_{i=1}^m \bigg(\frac{\eta}{\varepsilon + v^\ite_i} - \frac1 L \bigg)\norm{\uu^{\ite+1}_i - \uu^{\ite}_i}^2 \\
        &=\sum_{i\in I} \bigg(\frac{\eta}{\varepsilon + v^\ite_i} - \frac1 L \bigg)\norm{\uu^{\ite+1}_i - \uu^{\iter}_i}^2 + \sum_{i\in [m]\setminus I} \bigg(\frac{\eta}{\varepsilon + v^\ite_i} - \frac1 L \bigg)\norm{\uu^{\ite+1}_i - \uu^{\ite}_i}^2.
\end{align*}
Thus, for every $\iter\in \N$, with $\iter\geq \bar{\iter}$
\begin{align}
\label{eq:20260124c}
\nonumber\frac{1}{2L}\sum_{\ite= \bar{\iter}}^{\iter}\sum_{i\in I} \norm{\uu^{\ite+1}_i - \uu^{\ite}_i}^2 &\leq
\sum_{\ite= \bar{\iter}}^{\iter}\sum_{i\in I} \bigg(\frac1 L - \frac{\eta}{\varepsilon + v^\ite_i} \bigg)\norm{\uu^{\ite+1}_i - \uu^{\ite}_i}^2\\
\nonumber &\leq \sum_{\ite = \bar{\iter}}^{\iter} 2(F(\bx^{\ite}) - F(\bx^{\ite+1}))  
+ \sum_{\ite=\bar{\iter}}^{\iter}\sum_{i\in [m]\setminus I} \bigg(\frac{\eta}{\varepsilon + v^\ite_i} - \frac1 L \bigg)\norm{\uu^{\ite+1}_i - \uu^{\ite}_i}^2\\
\nonumber&\leq 2(F(\bx^{\bar{\iter}}) - F_*)  
+ \sum_{\ite=\bar{\iter}}^{\iter}\sum_{i\in [m]\setminus I} \frac{\eta}{\varepsilon + v^\ite_i} \norm{\uu^{\ite+1}_i - \uu^{\ite}_i}^2\\
\nonumber&\leq 2(F(\bx^{\bar{\iter}}) - F_*)  
+ \max_{1\leq i\leq m} \frac{\eta}{\varepsilon + v^{\bar{\iter}}_i} \sum_{\ite=\bar{\iter}}^{\iter} \sum_{i\in [m]\setminus I} \norm{\uu^{\ite+1}_i - \uu^{\ite}_i}^2\\
&\leq 2(F(\bx^{\bar{\iter}}) - F_*)  
+ \max_{1\leq i\leq m} \frac{\eta}{\varepsilon + v^{\bar{\iter}}_i} m(2\eta L- \varepsilon)^2,
\end{align}
where in the last inequality we used \eqref{eq:20260124b}. Thus, combining \eqref{eq:20260124c}
and \eqref{eq:20260124b} we get, for every integer $\iter\geq \bar{\iter}$,
\begin{align*}
\sum_{\ite=\bar{\iter}}^{\iter} \norm{\buu^{\ite+1}-\buu^{\ite}}^2 &= 
\sum_{\ite=\bar{\iter}}^{\iter} \sum_{i\in I} \norm{\uu_i^{\ite+1}-\uu_i^{\ite}}^2 
+\sum_{\ite=\bar{\iter}}^{\iter} \sum_{i\in [m]\setminus I} \norm{\uu_i^{\ite+1}-\uu_i^{\ite}}^2\\
&\leq 4L (F(\bx^{\bar{\iter}}) - F_*)  
+ 2L \max_{1\leq i\leq m} \frac{\eta}{\varepsilon + v^{\bar{\iter}}_i} m(2\eta L- \varepsilon)^2
+ m(2\eta L- \varepsilon)^2
\end{align*}
and the statement follows.
\end{proof}

%\begin{customcor}{\ref{cor:20260126a}}
%Let $(\x^\iter)_{\iter\in \N}$ be generated by Algorithm~\ref{adadiff}.
% with $\varepsilon>0$. 
% Then,
%for every $\iter\in \N$, $\W_\iter\succeq \varepsilon \Id$ and $\sup_{\iter\in \N}\norm{\W_\iter}<+\infty$.
%\end{customcor}
\begin{corollary}
\label{cor:20260126a}
Let $(\bx^\iter)_{\iter\in \N}$ be generated by Algorithm~\ref{adadiff}.
% with $\varepsilon>0$. 
 Then,
for all $\iter\in \N$, $\W_\iter\succeq \varepsilon \Id$ and $\sup_{\iter\in \N}\norm{\W_\iter}<+\infty$.
\end{corollary}
\begin{proof}
Recalling \eqref{AdaDiff} it is clear that, for every $\iter\in \N$ and $i\in\ [m]$, $\w_i^\iter\geq \varepsilon$, so that
$\W_\iter \succeq \varepsilon \Id$. Moreover,
\begin{equation*}
(\forall\,\iter\in \N)(\forall\,i\in [m])\quad w^\iter_i = \varepsilon + \sqrt{\sum_{\ite=1}^\iter \norm{\uu_i^\iter- \uu_i^{\iter-1}}^2} \leq \varepsilon + \sqrt{\sum_{\ite=1}^\iter \norm{\buu^\iter- \buu^{\iter-1}}^2},
\end{equation*}
so that $\sup_{\iter\in\ \N}\norm{\W^\iter}=\sup_{\iter\in\ \N}(\max_{1\leq i\leq m} \w_i^\iter) \leq \varepsilon + (\sum_{\iter=1}^{+\infty} \norm{\buu^\iter- \buu^{\iter-1}}^2)^{1/2}<+\infty$.
\end{proof}

\begin{definition}[{Definition~3.1(ii) in \cite{combettes2013variable}}]
\label{def:fejer}
Let $(V_\iter)_{\iter\in \N}$ be a sequence of positive self-adjoint bounded linear operators on the Hilbert space $\HH$,
such that, for every $\iter\in \N$, $V_n \succeq \alpha \Id$, for some $\alpha>0$.
Let $\norm{\cdot}_{V_n}$ be the Hilbert norms defined by the $V_\iter$'s
%${(\abs{\cdot}_\iter)}_{\iter \in \N}$ be a sequence of Hilbert norms on $\HH$
%such that, for some  $\nu>0$, $\nu \lVert \cdot\rVert^2 \leq \lvert \cdot\rvert ^2_\iter$
%for every $\iter \in \N$.
and let $S \subset \HH$ be a nonempty set.
A sequence ${(z^\iter)}_{\iter \in \N}$ in $\HH$ is a 
\emph{quasi-Fej\'er  sequence with respect to $S$
relative to ${(V_\iter)}_{\iter \in \N}$}  if
there exist summable sequences ${(\alpha_\iter)}_{\iter \in \N}$
and ${(\chi_\iter)}_{\iter \in \N}$ in $\R_{++}$ such that
\begin{equation}
\label{eq:fejerprop}
(\forall\, z \in S)(\forall\, \iter \in \N)\qquad \norm{z^{\iter+1} - z}_{V_{\iter+1}}^2 
\leq (1+\chi_\iter) \norm{z^\iter - z}_{V_\iter}^2 + \alpha_\iter.
\end{equation}
\end{definition}

\begin{fact}[Lemma~2.3, Proposition~3.2, and Theorem~3.3 in \cite{combettes2013variable}]
\label{fact:fejer}
Under the assumptions of Definition~\ref{def:fejer} suppose additionally that $\sup_{\iter\in \N} \norm{V_n}<+\infty$. Consider the two statements
\setlist[enumerate]{itemsep=0mm}
\begin{enumerate}[label={\rm (\alph*)}]
\vspace{-2ex}
\item\label{fact:fejer_a} either $(\forall\,\iter\in \N)\ V_{\iter+1} \succeq V_{\iter}$ or $(\forall\,\iter\in \N)\ V_{\iter+1} \preceq V_{\iter}$.
\item\label{fact:fejer_b} There exists a positive self-adjoint bounded linear operator $V\colon \HH\to \HH$ such that, for every $z\in \HH$, $V_\iter z \to Vz$.
\end{enumerate}
\vspace{-2ex}
Then \ref{fact:fejer_a} $\ \Rightarrow\ $ \ref{fact:fejer_b}. Moreover, suppose that $S\subset\HH$
is a nonempty set and let $(z^\iter)_{\iter\in \N}$ be a quasi-Fej\'er sequence in $\HH$ with respect to $S$
relative to $(V_\iter)_{\iter\in \N}$. Then if \ref{fact:fejer_b} holds, we have
\vspace{-2ex}
\begin{enumerate}[label={\rm (\roman*)}]
\item $(z^\iter)_{\iter\in \N}$ is bounded and, for every $z\in S$, $(\norm{z^\iter-z}_{V_\iter})_{\iter\in \N}$ is convergent
\item $(z^\iter)_{\iter\in \N}$ is weakly convergent to a point in $S$ iff every weak sequential cluster point of $(z^\iter)_{\iter\in \N}$ belongs to $S$.
\end{enumerate}
\end{fact}

\begin{customprop}{\ref{prop:Fejer}}
Under Assumptions~\ref{ass:0} and \ref{ass:2},
the sequence $(\bx^\iter)_{\iter\in \N}$ generated by Algorithm~\ref{adadiff} is a quasi-Fej\'er sequence with respect to $\argmin F$ relative to $(\W_\iter)_{\iter\in \N}$. More explicitly the following hold
\begin{equation*}
(\forall\,\iter\in \N)(\forall\,\bx\in \argmin F)\quad \norm{\bx^{\iter+1}-\bx}^2_{\W_{\iter+1}} \leq (1+\chi_\iter) \norm{\bx^{\iter}-\bx}^2_{\W_{\iter}} + \alpha_\iter
\end{equation*}
where
\begin{equation*}
\chi_\iter :=\max_{1\leq i\leq m}\frac{\varepsilon + v_i^{\iter}}{\varepsilon + v_i^{\iter-1}}- 1
\quad\text{and}\quad \alpha_\iter := \eta_\iter^2 \norm{\buu^{\iter+1} - \buu^{\iter}}_{\W_\iter^{-1}}^2
\end{equation*}
and $(\chi_\iter)_{\iter\in \N}$ and $(\alpha_\iter)_{\iter\in \N}$ are summable.
\end{customprop}
\begin{proof}
Let $\bx\in \argmin F$. Then it follows from inequality \eqref{lemma1_bound} in Lemma~\ref{lemma1} that
\begin{equation*}
(\forall\,\iter\in \N)\quad
\norm{\bx^{\iter+1} - \bx}_{\W_\iter}^2\leq \norm{\bx^\iter - \bx}_{\W_\iter}^2+ \eta_\iter^2 \norm{\buu^{\iter+1} - \buu^{\iter}}_{\W_\iter^{-1}}^2.
\end{equation*}
On the other hand
\begin{align*}
\norm{\bx^\iter - \bx}_{\W_\iter}^2 &= \sum_{i=1}^m (\varepsilon + v_i^\iter) \norm{\x_i^\iter- \x_i}^2
\leq \max_{1\leq i\leq m}\frac{\varepsilon + v_i^{\iter}}{\varepsilon + v_i^{\iter-1}}\sum_{i=1}^m (\varepsilon + v_i^\iter) \norm{\x_i^{\iter-1}- \x_i}^2\\[1ex]
& = \max_{1\leq i\leq m}\frac{\varepsilon + v_i^{\iter}}{\varepsilon + v_i^{\iter-1}}\norm{\bx^\iter - \bx}_{\W_{\iter-1}}^2 = (1+\chi_\iter)\norm{\bx^\iter - \bx}_{\W_{\iter-1}}^2,
\end{align*}
where, for every $\iter\in \N$, 
\begin{equation*}
0\leq \chi_\iter := \max_{1\leq i\leq m}\frac{\varepsilon + v_i^{\iter}}{\varepsilon + v_i^{\iter-1}}- 1
= \max_{1\leq i\leq m} \frac{v_i^{\iter} - v_i^{\iter-1}}{\varepsilon + v_i^{\iter-1}} \leq \max_{1\leq i\leq m} \frac{v_i^{\iter} - v_i^{\iter-1}}{\varepsilon} \leq \frac{1}{\varepsilon}\sum_{i=1}^m (v_i^{\iter} - v_i^{\iter-1}),
\end{equation*}
so that, 
\begin{equation*}
\sum_{\ite=1}^{\iter} \chi_\ite \leq \frac{1}{\varepsilon}\sum_{\ite=1}^{\iter} \sum_{i=1}^m(v_i^{\ite} - v_i^{\ite-1}) = \frac{1}{\varepsilon}\sum_{i=1}^m \sum_{\ite=1}^{\iter} (v_i^{\ite} - v_i^{\ite-1})\leq \frac{1}{\varepsilon} \sum_{i=1}^m v_i^\iter = \frac{1}{\varepsilon}\sum_{i=1}^{m} \sqrt{\sum_{\ite=1}^{\iter} \norm{\uu_i^{\ite}-\uu_i^{\ite-1}}^2}
\end{equation*}
and hence
$\sum_{\iter=1}^{+\infty} \chi_\iter \leq \varepsilon^{-1}m (\sum_{\iter=1}^{+\infty} \norm{\buu^{\iter}-\buu^{\iter-1}}^2)^{1/2}<+\infty$, by Proposition~\ref{proposition5}. Moreover,
\begin{equation*}
\alpha_\iter:=\eta_\iter^2 \norm{\buu^{\iter+1} - \buu^{\iter}}_{\W_\iter^{-1}}^2 = \eta^2 \norm{\buu^{\iter+1} - \buu^{\iter}}_{\W_\iter^{-1}}^2 \leq \frac{\eta^2}{\varepsilon} \norm{\buu^{\iter+1} - \buu^{\iter}}^2
\end{equation*}
and the statement follows again from Proposition~\ref{proposition5}.
\end{proof}

\begin{customprop}{\ref{prop:clusterpoints}}
Every weak sequential cluster point of $(\bx^{\iter})_{\iter\in \N}$ belongs to $\argmin F$.
\end{customprop}
\begin{proof}
First of all we show that there exists $(\bm{u}^{\iter})_{\iter\in \N}$ such that, for every $n\in \N$, $\bm{u}^\iter\in \partial \reg(\bx^\iter)$ and
 $\bm{u}^\iter\to 0$.
Let $\iter\in \N$.
As we saw in the proof of Lemma~\ref{lem:basicPGstep}, we have
\begin{equation*}
\frac{1}{\eta}\W_\iter(\bx^{\iter}-\bx^{\iter+1})- \buu^\iter\in \partial \reg(\bx^{\iter+1})
\end{equation*}
Therefore,
\begin{equation*}
\bm{u}^{\iter+1}:= \frac{1}{\eta}\W_\iter(\bx^{\iter}-\bx^{\iter+1}) + \buu^{\iter+1} - \buu^\iter \in \partial f(\bx^{\iter+1})+\partial \varphi(\bx^{\iter+1}) = \partial F(\bx^{\iter+1})
\end{equation*}
Now we observe that, setting $\beta= \sup_{\iter\in \N}\norm{\W_\iter}$ (recall Corollary~\ref{cor:20260126a}), we have
\begin{equation}
\label{eq:20260126b}
\norm{\bm{u}^{\iter+1}} \leq \frac{\beta}{\eta}\norm{\bx^{\iter+1}-\bx^\iter} + \norm{\buu^{\iter+1} - \buu^\iter}.
\end{equation}
Moreover, inequality \eqref{lemma1_bound} in Lemma~\ref{lemma1} with $\bx= \bx^\iter$ yields
\begin{equation*}
 \norm{\bx^{\iter+1} - \bx^\iter}_{\W_\iter}^2\leq 
2\eta(F(\bx^\iter)- F(\bx^{\iter+1}))  
+ \eta^2 \norm{\buu^{\iter+1} - \buu^{\iter}}_{\W_\iter^{-1}}^2
\end{equation*}
and since $\W_\iter \succeq \varepsilon \Id$ (by Corollary~\ref{cor:20260126a}) we have
\begin{equation*}
 \varepsilon \norm{\bx^{\iter+1} - \bx^\iter}^2\leq 
2\eta(F(\bx^\iter)- F(\bx^{\iter+1}))  
+ \frac{\eta^2}{\varepsilon} \norm{\buu^{\iter+1} - \buu^{\iter}}^2.
\end{equation*}
 In the end it follows from the above inequality and \eqref{eq:20260126b} that
 \begin{equation*}
 \norm{\bm{u}^{\iter+1}} \leq \frac{2\beta}{\varepsilon} (F(\bx^\iter)- F(\bx^{\iter+1}))  +
\bigg( \frac{\eta^2}{\varepsilon^2} + 1\bigg)
\norm{\buu^{\iter+1} - \buu^{\iter}}^2. 
 \end{equation*}
Thus, by Proposition~\ref{proposition5},
 \begin{equation*}
 \sum_{\iter=1}^{+\infty} \norm{\bm{u}^{\iter+1}} \leq \frac{2\beta}{\varepsilon}(F(\bx^1)- F_*)
 + \bigg( \frac{\eta^2}{\varepsilon^2} + 1\bigg)\sum_{\iter=1}^{+\infty} \norm{\buu^{\iter+1} - \buu^{\iter}}^2<+\infty
 \end{equation*}
 and hence $\bm{u}^{\iter}\to 0$. Now, let $(\bx^{\iter_k})_{k\in \N}$ be a subsequence of
 $(\bx^\iter)_{\iter\in \N}$ such that $\bx^{\iter_k} \rightharpoonup \bar{\bx}$. Then we have
 \begin{equation*}
 (\bx^{\iter_k}, \bm{u}^{\iter_k}) \in \mathrm{gra}(\partial F)\quad\text{and}\quad
 \bx^{\iter_k} \rightharpoonup \bar{\bx},  \bm{u}^{\iter_k} \to 0.
 \end{equation*}
 Since $F\in \Gamma_0(\bHH)$, the graph of $\partial F$ is weakly-strongly closed, meaning closed in the product of the topological spaces $\bHH_w\times \bHH$, where $\bHH_w$ is endowed with the weak topology and $\bHH$ is endowed with the strong topology.  Therefore $(\bar{\bx},0)\in \mathrm{gra}(\partial F)$ and hence $0\in \partial F(\bar{\bx})$. The statement follows.
\end{proof}
\section{Details on the Experiments and Additional Plots}\label{B}
\textbf{Details on the experiments.}
In the robustness experiments (first row of Figures \ref{fig:robustness_nonsmooth} and \ref{fig:robustness_smooth}), the horizontal axis corresponds to a grid of 200 candidate values of the stepsize parameter $\eta$.
In  Figures \ref{fig:loss-steps-splice}, \ref{fig:loss-steps-hinge}, \ref{fig:loss-steps-LAD}, \ref{fig:loss-steps-news20}, \ref{fig:loss-steps-2moons}, we report the optimality gap and corresponding stepsize (average over the coordinates) for three representative choices of $\eta.$ The central plot corresponds to a reference parameter $\eta$ obtained by averaging the best-performing values of $\eta$ (in terms of average performance) for AdaGrad and AdaGrad-Diff across the 10 random seeds.  The remaining plots use parameters obtained by scaling this reference value by factors of $10^{-1}$ and $10^{1}$, respectively, providing a simple heuristic for probing different operating regimes.\\

\noindent
When plotting the objective gap $F(\bx^\iter) - F_\star$, in general, the true optimal value $F_\star$ is unknown. We approximate it by keeping track of the smallest objective value observed across all methods, choices for $\eta$ in the grid considered, and iterations, over the 10 seeds. Additionally, the method with the fastest decay is run for 10 times the original budget to refine this estimate. All objective gaps are reported relative to this empirical $F_\star$.\\

\noindent
\textbf{Additional plots.} For completeness, we provide the same plots as those shown in the main body for all problem instances that were omitted there.
\begin{figure*}[ht]
    %\vskip 0.2in
    \centering
    \begin{subfigure}[t]{\textwidth}
        \centering
        \includegraphics[width=\linewidth]{./losses_hinge.png}
        \caption{Optimality gap.}
        \label{fig:losses-hinge}
    \end{subfigure}
    \hfill
    \begin{subfigure}[t]{\textwidth}
        \centering
        \includegraphics[width=\linewidth]{./stepsizes_hinge.png}
        \caption{Stepsizes evolution.}
        \label{fig:stepsizes-hinge}
    \end{subfigure}
    \caption{Comparison of AdaGrad and AdaGrad-Diff for the Hinge Loss on the synthetic dataset.
    (a) Optimality gaps  across three stepsizes parameter settings
    ($\eta = 0.0063,\, 0.063,\, 0.63$).
    (b) Stepsize evolution across different choices for the parameter $\eta$. The plots show the average and standard deviation over 10 initializations of the algorithms.}
    \label{fig:loss-steps-hinge}
\end{figure*}
\begin{figure*}[ht]
    %\vskip 0.2in
    \centering
    \begin{subfigure}[t]{\textwidth}
        \centering
        \includegraphics[width=\linewidth]{./losses_LAD.png}
        \caption{Optimality gap.}
        \label{fig:losses-splice}
    \end{subfigure}
    \hfill
    \begin{subfigure}[t]{\textwidth}
        \centering
        \includegraphics[width=\linewidth]{./stepsizes_LAD.png}
        \caption{Stepsizes evolution.}
        \label{fig:stepsizes-LAD}
    \end{subfigure}
    \caption{Comparison of AdaGrad and AdaGrad-Diff for LAD Regression on the synthetic dataset.
    (a) Optimality gaps  across three stepsizes parameter settings
    ($\eta = 0.0042,\, 0.042,\, 0.42$).
    (b) Stepsizes evolution across different choices for the parameter $\eta$. The plots show the average and standard deviation over 10 initializations of the algorithms.}
    \label{fig:loss-steps-LAD}
\end{figure*}
\begin{figure*}[ht]
    %\vskip 0.2in
    \centering
    \begin{subfigure}[t]{\textwidth}
        \centering
        \includegraphics[width=\linewidth]{./losses_news20.png}
        \caption{Optimality gap.}
        \label{fig:losses-news20}
    \end{subfigure}
    \hfill
    \begin{subfigure}[t]{\textwidth}
        \centering
        \includegraphics[width=\linewidth]{./stepsizes_news20.png}
        \caption{Stepsizes evolution.}
        \label{fig:stepsizes-news20}
    \end{subfigure}
    \caption{Comparison of AdaGrad and AdaGrad-Diff for Logistic Regression on the \texttt{news20.binary} dataset.
    (a) Optimality gaps  across three stepsizes parameter settings
    ($\eta = 0.0863,\, 0.863,\, 8.63$).
    (b) Stepsize evolution across different choices for the parameter $\eta$. The plots show the average and standard deviation over 10 initializations of the algorithms.}
    \label{fig:loss-steps-news20}
\end{figure*}
\begin{figure*}[ht]
    %\vskip 0.2in
    \centering
    \begin{subfigure}[t]{\textwidth}
        \centering
        \includegraphics[width=\linewidth]{./losses_2moons.png}
        \caption{Optimality gap.}
    \end{subfigure}
    \hfill
    \begin{subfigure}[t]{\textwidth}
        \centering
        \includegraphics[width=\linewidth]{./stepsizes_2moons.png}
        \caption{Stepsizes evolution.}
    \end{subfigure}
    \caption{Comparison of AdaGrad and AdaGrad-Diff for SVM Classification on the \texttt{2moons} dataset.
    (a) Optimality gaps  across three stepsizes parameter settings
    ($\eta = 0.0002,\, 0.002,\, 0.02$).
    (b) Stepsize evolution across different choices for the parameter $\eta$. The plots show the average and standard deviation over 10 initializations of the algorithms.}
    \label{fig:loss-steps-2moons}
\end{figure*}

%You can have as much text here as you want. The main body must be at most $8$ pages long. For the final version, one more page can be added. If you want, you can use an appendix like this one.

%The $\mathtt{\backslash onecolumn}$ command above can be kept in place if you prefer a one-column appendix, or can be removed if you prefer a two-column appendix.  Apart from this possible change, the style (font size, spacing, margins, page numbering, etc.) should be kept the same as the main body.
%%%%%%%%%%%%%%%%%%%%%%%%%%%%%%%%%%%%%%%%%%%%%%%%%%%%%%%%%%%%%%%%%%%%%%%%%%%%%%%
%%%%%%%%%%%%%%%%%%%%%%%%%%%%%%%%%%%%%%%%%%%%%%%%%%%%%%%%%%%%%%%%%%%%%%%%%%%%%%%

\end{document}

% This document was modified from the file originally made available by
% Pat Langley and Andrea Danyluk for ICML-2K. This version was created
% by Iain Murray in 2018, and modified by Alexandre Bouchard in
% 2019 and 2021 and by Csaba Szepesvari, Gang Niu and Sivan Sabato in 2022.
% Modified again in 2023 and 2024 by Sivan Sabato and Jonathan Scarlett.
% Previous contributors include Dan Roy, Lise Getoor and Tobias
% Scheffer, which was slightly modified from the 2010 version by
% Thorsten Joachims & Johannes Fuernkranz, slightly modified from the
% 2009 version by Kiri Wagstaff and Sam Roweis's 2008 version, which is
% slightly modified from Prasad Tadepalli's 2007 version which is a
% lightly changed version of the previous year's version by Andrew
% Moore, which was in turn edited from those of Kristian Kersting and
% Codrina Lauth. Alex Smola contributed to the algorithmic style files.
